\documentclass[aspectratio=169]{beamer}
\usetheme{UFS}

\usepackage[utf8]{inputenc}
\usepackage[portuguese]{babel}
\usepackage{amsmath,amssymb,graphicx,booktabs,xcolor}
\usepackage{tikz}
\usetikzlibrary{arrows.meta,shapes,shapes.geometric,positioning,calc,fit,
  backgrounds,decorations.pathreplacing,matrix,patterns}
\usepackage{listings}

% UFS colors (explicit para uso no TikZ)
\definecolor{UFSBlue}{RGB}{3,78,162}
\definecolor{UFSYellow}{RGB}{253,185,19}
\definecolor{UFSRed}{RGB}{237,28,36}
\definecolor{UFSGray}{RGB}{100,100,100}
\definecolor{UFSLightGray}{RGB}{230,230,230}

% Listing style Java
\definecolor{kwcolor}{HTML}{0000FF}
\definecolor{strcolor}{HTML}{A31515}
\definecolor{cmtcolor}{HTML}{008000}
\definecolor{numcolor}{HTML}{888888}

\lstdefinestyle{java}{
  language=Java,
  basicstyle=\ttfamily\scriptsize,
  keywordstyle=\color{kwcolor}\bfseries,
  stringstyle=\color{strcolor},
  commentstyle=\color{cmtcolor}\itshape,
  numberstyle=\tiny\color{numcolor},
  numbers=left, numbersep=6pt,
  xleftmargin=14pt,
  breaklines=true, tabsize=2,
  showstringspaces=false,
  frame=single, rulecolor=\color{UFSBlue!30},
  morekeywords={var,record,sealed,permits,yield,instanceof},
}
\lstset{style=java}

% Estilos TikZ reutilizáveis
\tikzset{
  ufnode/.style={circle,draw=UFSBlue,fill=UFSBlue!15,thick,
    font=\small\bfseries,minimum size=0.65cm,inner sep=1pt},
  ufroot/.style={circle,draw=UFSRed,fill=UFSRed!20,very thick,
    font=\small\bfseries,minimum size=0.65cm,inner sep=1pt},
  ufedge/.style={-{Stealth[scale=0.8]},thick,UFSGray!80},
  llbox/.style={rectangle,draw=UFSBlue,fill=UFSBlue!10,
    minimum width=0.85cm,minimum height=0.58cm,font=\scriptsize,inner sep=2pt},
  llrep/.style={rectangle,draw=UFSRed,fill=UFSRed!10,very thick,
    minimum width=0.85cm,minimum height=0.58cm,font=\scriptsize\bfseries,inner sep=2pt},
  step/.style={rectangle,rounded corners=3pt,draw=UFSGray,fill=UFSLightGray,
    minimum width=1.5cm,minimum height=0.55cm,font=\scriptsize,align=center},
}

\title{Conjuntos Disjuntos}
\subtitle{Union-Find, Path Compression e Representação por Floresta}
\author{Prof.\ Andre Yoshiaki Kashiwabara}
\institute{DCOMP -- Universidade Federal de Sergipe\\[2pt]
\texttt{andreyoshiaki@dcomp.ufs.br}}
\date{}

\begin{document}

%% ================================================================
%%  ABERTURA
%% ================================================================

\begin{frame}[plain]%% 1
  \titlepage
\end{frame}


\begin{frame}{O que você vai aprender}%% 3
  \begin{center}
  \begin{tikzpicture}[
    box/.style={rectangle,rounded corners=4pt,draw=UFSBlue,fill=UFSBlue!10,
      minimum width=3.0cm,minimum height=0.68cm,font=\small\bfseries,align=center},
    goal/.style={rectangle,rounded corners=4pt,draw=UFSRed,fill=UFSRed!10,
      minimum width=3.0cm,minimum height=0.68cm,font=\small\bfseries,align=center},
    ar/.style={-{Stealth},thick,UFSGray}
  ]
    %% Linha 1 – topo
    \node[box] (adt)    at (0.5,3.0)   {Problema e ADT};
    %% Linha 2 – duas alternativas
    \node[box] (ll)     at (-3.0,2.0)  {Lista Encadeada};
    \node[box] (forest) at (2.0,2.0)   {Floresta};
    %% Linha 3 – heurísticas (lista: peso; floresta: rank, PC)
    \node[box] (peso)   at (-3.8,1.0)  {Union por Peso};
    \node[box] (rank)   at (-0.5,1.0)  {Union by Rank};
    \node[box] (pc)     at (3.5,1.0)   {Path Compression};
    %% Linha 4 – resultado e implementação
    \node[goal] (ana)   at (1.0,0.0)   {Análise: $O(m\,\alpha(n))$};
    \node[goal] (java)  at (4.6,0.0)   {Java Moderno};
    %% Nota sobre alpha(n) abaixo do nó de análise
    \node[font=\scriptsize,UFSGray,align=center,text width=4.0cm]
      at (1.0,-0.75)
      {($\alpha$: inversa de Ackermann — veremos adiante)};
    %% Setas
    \draw[ar] (adt)--(ll);
    \draw[ar] (adt)--(forest);
    \draw[ar] (ll)--(peso);
    \draw[ar] (peso)--(ana);
    \draw[ar] (forest)--(rank);
    \draw[ar] (forest)--(pc);
    \draw[ar] (rank)--(ana);
    \draw[ar] (pc)--(ana);
    \draw[ar] (ana)--(java);
  \end{tikzpicture}
  \end{center}
\end{frame}

%% ================================================================
%%  SEÇÃO 1 — MOTIVAÇÃO
%% ================================================================
\section{Motivação}

\begin{frame}{Conectividade dinâmica}%% 4
  \begin{columns}[T]
    \begin{column}{0.55\textwidth}
      \begin{block}{Problema}
        Dados $n$ objetos, processar duas operações:
        \begin{itemize}
          \item \textbf{Union}$(p,q)$: registrar que $p$ e $q$ estão conectados
          \item \textbf{Connected}$(p,q)$: $p$ e $q$ estão no mesmo componente?
        \end{itemize}
      \end{block}
      \vspace{0.25cm}
      Conectividade é transitiva: se $p$--$q$ e $q$--$r$, então $p$--$r$.\\
      Matematicamente: \textbf{relação de equivalência}.
    \end{column}
    \begin{column}{0.42\textwidth}
      \begin{tikzpicture}[
        n/.style={circle,draw=UFSBlue,fill=UFSBlue!15,minimum size=0.58cm,font=\scriptsize\bfseries}
      ]
        \foreach \i/\x/\y in {0/0/2.6,1/1.5/2.6,2/3/2.6,
                               3/0/1.2, 4/1.5/1.2,5/3/1.2}
          \node[n] (\i) at (\x,\y) {\i};
        \draw[thick,UFSBlue] (0)--(3);
        \draw[thick,UFSBlue] (1)--(4);
        \draw[thick,UFSBlue] (3)--(4);
        \draw[thick,UFSRed,dashed] (2)--(5);
        \node[font=\tiny,UFSGray,below=0.1cm of 4]
          {$\{0,1,3,4\}$\quad$\{2,5\}$};
      \end{tikzpicture}
    \end{column}
  \end{columns}
\end{frame}

\begin{frame}{Por que isso importa?}%% 5
  \begin{columns}[T]
    \begin{column}{0.48\textwidth}
      \begin{block}{Algoritmos de grafos}
        \begin{itemize}
          \item Kruskal — MST em $O(m\,\alpha(n))$
          \item Componentes conexas on-line
          \item Detecção de ciclos
          \item Menor caminho com pesos (Borůvka)
        \end{itemize}
      \end{block}
    \end{column}
    \begin{column}{0.48\textwidth}
      \begin{block}{Aplicações em sistemas}
        \begin{itemize}
          \item Unificação de tipos (compiladores)
          \item Segmentação de imagens
          \item Percolação (Monte Carlo)
          \item Particionamento de memória (GC)
        \end{itemize}
      \end{block}
    \end{column}
  \end{columns}
\end{frame}

\begin{frame}{Interface: Conjunto Disjunto (CLRS, Cap.\ 19)}%% 6
  \begin{block}{Três operações fundamentais}
    \begin{description}
      \item[\textsc{Make-Set}$(x)$] Cria um conjunto $\{x\}$; $x$ é seu próprio representante.
      \item[\textsc{Find-Set}$(x)$] Retorna o \emph{representante} do conjunto que contém $x$.
      \item[\textsc{Union}$(x,y)$] Funde os conjuntos de $x$ e $y$ em um único conjunto.
    \end{description}
  \end{block}
  \vspace{0.4cm}
  \begin{alertblock}{Invariante fundamental}
    $\textsc{Find-Set}(x) = \textsc{Find-Set}(y) \iff x$ e $y$ pertencem ao mesmo conjunto.
  \end{alertblock}
\end{frame}

\begin{frame}[shrink=8]{Exemplo: as três operações em ação}%% 6b
  \begin{center}
  \begin{tikzpicture}[
    nd/.style={circle,draw=UFSBlue,fill=UFSBlue!10,minimum size=0.6cm,font=\scriptsize\bfseries},
    rt/.style={circle,draw=UFSRed,fill=UFSRed!15,very thick,minimum size=0.6cm,font=\scriptsize\bfseries},
    ar/.style={-{Stealth[scale=0.7]},thick,UFSGray},
    lbl/.style={font=\scriptsize,UFSBlue},
    op/.style={font=\scriptsize\ttfamily,UFSRed}
  ]
    %% Linha 0: Make-Set(1..4)  — labels fundidos para evitar sobreposição
    \node[lbl,align=right] at (-1.5,3.5)
      {\textsc{Make-Set}\\(1..4):};
    \foreach \i in {1,...,4}
      \node[rt] at ({(\i-1)*1.2},3.5) {\i};

    %% Linha 1: Union(1,2) — root em y=2.4, filho em y=1.6
    \node[op] at (-1.5,2.4) {Union(1,2):};
    \node[rt,label={[font=\tiny]above:rep}] (r1) at (0.6,2.4) {1};
    \node[nd] (n2) at (0.6,1.6) {2};          %% gap para próxima linha: 1.6-0.8=0.8cm>0.6
    \draw[ar] (n2)--(r1);
    \node[rt] at (2.4,2.4) {3};
    \node[rt] at (3.6,2.4) {4};

    %% Linha 2: Union(3,4) — root em y=0.8, filho em y=0.0
    \node[op] at (-1.5,0.8) {Union(3,4):};
    \node[rt] (r12) at (0.6,0.8) {1};
    \node[nd] (n12) at (0.6,0.0) {2};
    \draw[ar] (n12)--(r12);
    \node[rt] (r3) at (3.0,0.8) {3};
    \node[nd] (n4) at (3.0,0.0) {4};
    \draw[ar] (n4)--(r3);

    %% Linha 3: Find-Set(4) -> 3  — root em y=-1.2, filho em y=-2.0
    \node[op] at (-1.5,-1.2) {Find-Set(4):};
    \node[rt] (f3) at (3.0,-1.2) {3};
    \node[nd] (f4) at (3.0,-2.0) {4};
    \draw[ar] (f4)--(f3);
    \draw[-{Stealth},very thick,UFSYellow!80!black] (f4)--++(0.7,0)
      node[right,font=\scriptsize,UFSBlue] {retorna \textbf{3}};

    %% Linha 4: Union(1,3) — painel direito
    \node[op] at (5.5,2.4) {Union(1,3):};
    \node[rt,label={[font=\tiny]above:rep\,de\,tudo}] (ur) at (6.8,2.4) {1};
    \node[nd] (u2) at (6.0,1.5) {2};
    \node[nd] (u3) at (6.8,1.5) {3};
    \node[nd] (u4) at (7.6,0.6) {4};
    \draw[ar] (u2)--(ur); \draw[ar] (u3)--(ur); \draw[ar] (u4)--(u3);
    \node[font=\tiny,UFSGray] at (6.8,0.0) {$\{1,2,3,4\}$};
  \end{tikzpicture}
  \end{center}
\end{frame}

\begin{frame}{Cenário de uso típico}%% 7
  \begin{enumerate}
    \item \textbf{Fase de inicialização}: $n$ chamadas a \textsc{Make-Set} — uma por objeto.
    \item \textbf{Fase de processamento}: as $m-n$ operações restantes são \textsc{Union} (no máximo $n-1$) e \textsc{Find-Set}.
  \end{enumerate}
  \vspace{0.4cm}
  \begin{exampleblock}{Meta de desempenho}
    Processar $m$ operações em tempo total $O(m\,\alpha(n))$,\\
    onde $\alpha$ é a \textbf{inversa de Ackermann}: $\alpha(n)\leq 4$ para qualquer $n$ concebível.
  \end{exampleblock}
  \vspace{0.3cm}
  \begin{center}
    Na prática, é \textbf{praticamente linear}: $O(m\,\alpha(n))$ com $\alpha(n)\leq 4$.
  \end{center}
\end{frame}

\begin{frame}{Evolução das implementações}%% 8
  \begin{center}
  \renewcommand{\arraystretch}{1.4}
  \begin{tabular}{lll}
    \toprule
    \textbf{Implementação} & \textbf{Find-Set} & \textbf{Union} \\
    \midrule
    Lista encadeada ingênua        & $O(1)$         & $O(n)$ \\
    Lista + union por peso         & $O(1)$         & amort.\ $O(\lg n)$ \\
    Floresta ingênua               & $O(n)$         & $O(1)$ \\
    Floresta + union por rank      & $O(\lg n)$     & $O(1)$ \\
    Floresta + path compression    & amort.\ $O(\lg n)$ & $O(1)$ \\
    \textbf{Floresta + rank + PC}  & $\mathbf{O(\alpha(n))}$ & $\mathbf{O(1)}$ \\
    \bottomrule
  \end{tabular}
  \end{center}
\end{frame}

\begin{frame}{Recapitulando: Motivação}%% 9
  \begin{block}{O que aprendemos}
    \begin{itemize}
      \item Conjuntos disjuntos modelam relações de equivalência dinâmicas
      \item Três operações: \textsc{Make-Set}, \textsc{Find-Set}, \textsc{Union}
      \item Objetivo: custo total $O(m\,\alpha(n))$ — praticamente linear
      \item Duas heurísticas simples permitem atingir esse resultado
    \end{itemize}
  \end{block}
\end{frame}

%% ================================================================
%%  SEÇÃO 2 — LISTA ENCADEADA
%% ================================================================
\section{Representação por Lista Encadeada}

\begin{frame}{Estrutura: lista encadeada}%% 10
  \begin{columns}[T]
    \begin{column}{0.52\textwidth}
      Cada conjunto = uma \textbf{lista encadeada}:
      \begin{itemize}
        \item Cada nó tem campo \texttt{rep} → representante (head)
        \item A lista mantém ponteiros \texttt{head} e \texttt{tail}
        \item \textsc{Find-Set}$(x)$: retorna $x.\texttt{rep}$ em $O(1)$
        \item \textsc{Union}: concatena e atualiza todos os \texttt{rep}
      \end{itemize}
    \end{column}
    \begin{column}{0.45\textwidth}
      \begin{tikzpicture}[
        nd/.style={llbox},
        rp/.style={llrep},
        ar/.style={-{Stealth[scale=0.7]},thick,UFSGray}
      ]
        \node[rp] (h) at (0,2.2) {$c_1$};
        \node[nd] (a) at (1.2,2.2) {$c_2$};
        \node[nd] (b) at (2.4,2.2) {$c_3$};
        \draw[ar] (h)--(a); \draw[ar] (a)--(b);
        \draw[ar,bend right=40,UFSRed] (a) to (h);
        \draw[ar,bend right=40,UFSRed] (b) to (h);
        \node[font=\tiny,UFSGray] at (1.2,1.4)
          {ponteiros \texttt{rep} em vermelho};
        \node[font=\tiny] at (-0.1,2.6) {\textsf{head}};
        \node[font=\tiny] at (2.4,2.6) {\textsf{tail}};
      \end{tikzpicture}
    \end{column}
  \end{columns}
\end{frame}

\begin{frame}[shrink=5]{Make-Set e Find-Set na lista}%% 11
  \begin{columns}[T]
    \begin{column}{0.48\textwidth}
      \begin{block}{\textsc{Make-Set}$(x)$}
        $x.\texttt{next} \leftarrow \text{NIL}$\\
        $x.\texttt{rep}  \leftarrow x$\\
        $\text{novo conjunto}: \{\texttt{head}=x,\,\texttt{tail}=x\}$\\[4pt]
        \textbf{Custo}: $O(1)$
      \end{block}
    \end{column}
    \begin{column}{0.48\textwidth}
      \begin{block}{\textsc{Find-Set}$(x)$}
        \textbf{return} $x.\texttt{rep}$\\[4pt]
        \textbf{Custo}: $O(1)$ — acesso direto.
      \end{block}
    \end{column}
  \end{columns}
  \vspace{0.5cm}
  \begin{exampleblock}{Exemplo: Make-Set(1), Make-Set(2), Make-Set(3)}
    \begin{center}
    \begin{tikzpicture}[
      nd/.style={llrep},
      ar/.style={-{Stealth[scale=0.7]},thick,UFSRed}
    ]
      \node[nd] (x1) at (0,0) {1};
      \node[nd] (x2) at (2,0) {2};
      \node[nd] (x3) at (4,0) {3};
      \draw[ar] (x1) to[out=120,in=60,looseness=5] (x1);
      \draw[ar] (x2) to[out=120,in=60,looseness=5] (x2);
      \draw[ar] (x3) to[out=120,in=60,looseness=5] (x3);
      \node[font=\tiny,UFSGray] at (0,-0.6) {$\{1\}$};
      \node[font=\tiny,UFSGray] at (2,-0.6) {$\{2\}$};
      \node[font=\tiny,UFSGray] at (4,-0.6) {$\{3\}$};
    \end{tikzpicture}
    \end{center}
  \end{exampleblock}
\end{frame}

\begin{frame}{Union na lista: algoritmo}%% 12
  \begin{block}{\textsc{Union}$(x,y)$ — sem heurística}
    \begin{enumerate}
      \item $S_x \leftarrow \text{conjunto de } x$;\quad $S_y \leftarrow \text{conjunto de } y$
      \item $S_x.\texttt{tail}.\texttt{next} \leftarrow S_y.\texttt{head}$
      \item \textbf{Para cada} nó $z$ em $S_y$:\ $z.\texttt{rep} \leftarrow S_x.\texttt{head}$
      \item $S_x.\texttt{tail} \leftarrow S_y.\texttt{tail}$
    \end{enumerate}
  \end{block}
  \vspace{0.2cm}
  \begin{alertblock}{Custo}
    Passo 3 percorre toda $S_y$ $\Rightarrow$ $O(|S_y|)$.\\
    Sem heurística, sequência adversarial leva a $\Theta(n^2)$.
  \end{alertblock}
\end{frame}

\begin{frame}{Passo a passo: Union na lista}%% 13
  \begin{center}
  \begin{tikzpicture}[
    nd/.style={llbox},
    rp/.style={llrep},
    ar/.style={-{Stealth[scale=0.65]},thick,UFSGray}
  ]
    %% Estado inicial: S1={a,b,c}  S2={d,e}
    \node[font=\scriptsize\bfseries,UFSBlue] at (-0.7,3.0) {$S_1$:};
    \node[rp] (a) at (0,3.0) {$a$};
    \node[nd] (b) at (1.1,3.0) {$b$};
    \node[nd] (c) at (2.2,3.0) {$c$};
    \draw[ar] (a)--(b); \draw[ar] (b)--(c);
    \draw[ar,bend right=30,UFSRed] (b) to (a);
    \draw[ar,bend right=30,UFSRed] (c) to (a);

    \node[font=\scriptsize\bfseries,UFSRed] at (-0.7,1.8) {$S_2$:};
    \node[rp,fill=UFSRed!20] (d) at (0,1.8) {$d$};
    \node[nd] (e) at (1.1,1.8) {$e$};
    \draw[ar] (d)--(e);
    \draw[ar,bend right=30,UFSRed] (e) to (d);

    %% Seta de operação
    \node[font=\small] at (5.2,2.4) {\textsc{Union}$(a,d)$};
    \draw[-{Stealth},thick,UFSGray] (3.2,2.4)--(4.0,2.4);

    %% Resultado
    \node[font=\scriptsize\bfseries,UFSBlue] at (-0.7,0.4) {$S'$:};
    \node[rp] (r1) at (0,0.4) {$a$};
    \node[nd] (r2) at (1.1,0.4) {$b$};
    \node[nd] (r3) at (2.2,0.4) {$c$};
    \node[nd] (r4) at (3.3,0.4) {$d$};
    \node[nd] (r5) at (4.4,0.4) {$e$};
    \draw[ar] (r1)--(r2); \draw[ar] (r2)--(r3);
    \draw[ar] (r3)--(r4); \draw[ar] (r4)--(r5);
    \draw[ar,bend right=30,UFSRed] (r4) to (r1);
    \draw[ar,bend right=30,UFSRed] (r5) to (r1);
    \node[font=\tiny,UFSGray] at (2.2,-0.4)
      {ponteiros \texttt{rep} de $d$ e $e$ atualizados para $a$ — custo $O(|S_2|)$};
  \end{tikzpicture}
  \end{center}
\end{frame}

\begin{frame}{Pior caso sem heurística}%% 14
  \begin{alertblock}{Sequência adversarial}
    \textsc{Union}$(x_1,x_2)$, \textsc{Union}$(x_2,x_3)$, \ldots, \textsc{Union}$(x_{n-1},x_n)$\\
    sempre concatenando lista de 1 elemento na lista crescente:\\
    custo $= 1 + 2 + \cdots + (n-1) = \Theta(n^2)$.
  \end{alertblock}
  \vspace{0.3cm}
  \begin{center}
  \begin{tikzpicture}[
    nd/.style={llbox},
    rp/.style={llrep},
    ar/.style={-{Stealth[scale=0.65]},thick,UFSGray}
  ]
    \node[rp] (x1) at (0,0) {$x_1$};
    \node[nd] (x2) at (1.1,0) {$x_2$};
    \node[nd] (x3) at (2.2,0) {$x_3$};
    \node at (3.3,0) {$\cdots$};
    \node[nd] (xn) at (4.4,0) {$x_n$};
    \draw[ar] (x1)--(x2); \draw[ar] (x2)--(x3);
    \draw[ar] (x3)--(2.9,0); \draw[ar] (3.7,0)--(xn);
    \draw[ar,bend right=35,UFSRed] (x2) to (x1);
    \draw[ar,bend right=35,UFSRed] (x3) to (x1);
    \draw[ar,bend right=35,UFSRed] (xn) to (x1);
  \end{tikzpicture}
  \end{center}
\end{frame}

\begin{frame}{Heurística: Union por Peso}%% 15
  \begin{block}{Regra (CLRS, Seção 19.2)}
    Anexar sempre a lista \textbf{menor} ao final da lista \textbf{maior}.\\
    O representante guarda o \textbf{tamanho} do conjunto.
  \end{block}
  \vspace{0.3cm}
  \begin{columns}[T]
    \begin{column}{0.5\textwidth}
      \begin{exampleblock}{Lema chave}
        Cada nó tem seu ponteiro \texttt{rep} atualizado \textbf{no máximo} $\lfloor\lg n\rfloor$ vezes.\\[4pt]
        \textbf{Por quê?} A cada atualização, o conjunto ao qual $x$ pertence ao menos \textbf{dobra} de tamanho.
      \end{exampleblock}
    \end{column}
    \begin{column}{0.47\textwidth}
      \begin{block}{Teorema 19.1 (CLRS)}
        Com $m$ operações totais sobre $n$ elementos:\\[4pt]
        \textbf{Custo total}: $O(m + n\lg n)$
      \end{block}
    \end{column}
  \end{columns}
\end{frame}

\begin{frame}{Passo a passo: Union por Peso}%% 16
  \begin{center}
  \begin{tikzpicture}[
    nd/.style={llbox},
    rp/.style={llrep},
    ar/.style={-{Stealth[scale=0.65]},thick,UFSGray},
    lbl/.style={font=\tiny,UFSGray}
  ]
    %% Passo 1: {a,b,c} tamanho 3  {d} tamanho 1
    \node[lbl] at (-1.0,3.2) {Passo 1:};
    \node[rp] (p1a) at (0,3.2) {$a$\,(3)};
    \node[nd] (p1b) at (1.4,3.2) {$b$};
    \node[nd] (p1c) at (2.6,3.2) {$c$};
    \node[rp,fill=UFSRed!15] (p1d) at (4.0,3.2) {$d$\,(1)};
    \draw[ar] (p1a)--(p1b); \draw[ar] (p1b)--(p1c);

    %% Seta
    \draw[-{Stealth},thick,UFSBlue!50] (5.2,3.2)--(5.8,3.2);
    \node[lbl] at (7.6,3.2) {$|S_a|{>}|S_d|$: anexar $d$ em $a$};

    %% Resultado passo 1
    \node[lbl] at (-1.0,2.0) {Resultado:};
    \node[rp] (r1a) at (0,2.0) {$a$\,(4)};
    \node[nd] (r1b) at (1.4,2.0) {$b$};
    \node[nd] (r1c) at (2.6,2.0) {$c$};
    \node[nd] (r1d) at (3.8,2.0) {$d$};
    \draw[ar] (r1a)--(r1b); \draw[ar] (r1b)--(r1c); \draw[ar] (r1c)--(r1d);
    \draw[ar,bend right=25,UFSRed] (r1d) to (r1a);

    %% Passo 2: unir {a,b,c,d} com {e,f} tamanho 2
    \node[lbl] at (-1.0,0.8) {Passo 2:};
    \node[rp,fill=UFSRed!15] (p2e) at (0,0.8) {$e$\,(2)};
    \node[nd] (p2f) at (1.4,0.8) {$f$};
    \draw[ar] (p2e)--(p2f);
    \node[lbl] at (5.5,0.8) {$|S_a|{>}|S_e|$: anexar $e,f$ em $a$};
    \draw[-{Stealth},thick,UFSBlue!50] (2.8,0.8)--(3.4,0.8);
    \node[lbl] at (-1.0,-0.3) {Resultado:};
    \node[rp] (res) at (0.3,-0.3) {$a$\,(6)};
    \node[lbl] at (2.5,-0.3) {$b\!\to\!c\!\to\!d\!\to\!e\!\to\!f$\ (todos com \texttt{rep}$=a$)};
  \end{tikzpicture}
  \end{center}
\end{frame}

\begin{frame}{Recapitulando: Lista Encadeada}%% 17
  \begin{block}{O que aprendemos}
    \begin{itemize}
      \item \textsc{Find-Set} em $O(1)$; \textsc{Union} em $O(n)$ sem heurística
      \item Union por peso: cada nó é atualizado $\leq \lfloor\lg n\rfloor$ vezes
      \item Custo total: $O(m + n\lg n)$ com union por peso
      \item Ainda caro para $n$ grande: precisamos evitar percorrer a lista no Union
    \end{itemize}
  \end{block}
\end{frame}

%% ================================================================
%%  SEÇÃO 3 — FLORESTA
%% ================================================================
\section{Representação por Floresta}

\begin{frame}{Ideia: árvores enraizadas}%% 18
  \begin{columns}[T]
    \begin{column}{0.5\textwidth}
      \begin{block}{Representação por floresta}
        \begin{itemize}
          \item Cada conjunto = uma \textbf{árvore enraizada}
          \item Raiz = representante do conjunto
          \item Cada nó guarda apenas o ponteiro para o \textbf{pai}
          \item Raiz: $x.\texttt{parent} = x$
        \end{itemize}
      \end{block}
    \end{column}
    \begin{column}{0.47\textwidth}
      \begin{tikzpicture}
        \node[ufroot] (r1) at (0.8,2.8) {1};
        \node[ufnode] (n2) at (0,1.8) {2};
        \node[ufnode] (n3) at (1.6,1.8) {3};
        \node[ufnode] (n4) at (0,0.8) {4};
        \draw[ufedge] (n2)--(r1); \draw[ufedge] (n3)--(r1);
        \draw[ufedge] (n4)--(n2);
        \node[ufroot] (r5) at (3.2,2.8) {5};
        \node[ufnode] (n6) at (3.2,1.8) {6};
        \draw[ufedge] (n6)--(r5);
        \node[font=\tiny,UFSGray] at (2,0.2)
          {$S_1=\{1,2,3,4\}$\quad $S_2=\{5,6\}$};
      \end{tikzpicture}
    \end{column}
  \end{columns}
\end{frame}

\begin{frame}{Passo a passo: Make-Set na floresta}%% 19
  \begin{block}{\textsc{Make-Set}$(x)$}
    $x.\texttt{parent} \leftarrow x$\\
    $x.\texttt{rank}   \leftarrow 0$\quad (preparando para Union by Rank)\\[4pt]
    \textbf{Custo}: $O(1)$
  \end{block}
  \vspace{0.3cm}
  \begin{exampleblock}{Exemplo: Make-Set(1..6) — seis conjuntos singleton}
    \begin{center}
    \begin{tikzpicture}
      \foreach \i in {1,...,6}
        \node[ufroot] at ({(\i-1)*1.2},0) {\i};
      \node[font=\tiny,UFSGray] at (3,-0.7)
        {cada nó é raiz de si mesmo — rank 0};
    \end{tikzpicture}
    \end{center}
  \end{exampleblock}
\end{frame}

\begin{frame}{Passo a passo: Find-Set na floresta}%% 20
  \begin{columns}[T]
    \begin{column}{0.5\textwidth}
      \begin{block}{\textsc{Find-Set}$(x)$ — sem compressão}
        \textbf{while} $x.\texttt{parent} \neq x$:\\
        \quad $x \leftarrow x.\texttt{parent}$\\
        \textbf{return} $x$\\[4pt]
        \textbf{Custo}: $O(\text{altura da árvore})$
      \end{block}
    \end{column}
    \begin{column}{0.47\textwidth}
      \begin{tikzpicture}
        \node[ufroot] (r)  at (1.2,3.2) {1};
        \node[ufnode] (n2) at (0.4,2.2) {2};
        \node[ufnode] (n3) at (2.0,2.2) {3};
        \node[ufnode] (n4) at (0.4,1.2) {4};
        \node[ufnode] (n7) at (0.4,0.2) {7};
        \draw[ufedge] (n2)--(r); \draw[ufedge] (n3)--(r);
        \draw[ufedge] (n4)--(n2); \draw[ufedge] (n7)--(n4);
        % Highlight path from 7 to root
        \draw[-{Stealth[scale=0.9]},very thick,UFSYellow!80!black]
          (n7) to[bend right=25] (n4);
        \draw[-{Stealth[scale=0.9]},very thick,UFSYellow!80!black]
          (n4) to[bend right=25] (n2);
        \draw[-{Stealth[scale=0.9]},very thick,UFSYellow!80!black]
          (n2) to[bend right=25] (r);
        \node[font=\tiny,UFSGray] at (1.2,-0.5)
          {\textsc{Find-Set}(7): caminho amarelo até raiz 1};
      \end{tikzpicture}
    \end{column}
  \end{columns}
\end{frame}

\begin{frame}[shrink=5]{Passo a passo: Union na floresta}%% 21
  \begin{columns}[T]
    \begin{column}{0.5\textwidth}
      \begin{block}{\textsc{Union}$(x,y)$ — sem rank}
        $r_x \leftarrow \textsc{Find-Set}(x)$\\
        $r_y \leftarrow \textsc{Find-Set}(y)$\\
        \textbf{if} $r_x \neq r_y$:\\
        \quad $r_y.\texttt{parent} \leftarrow r_x$\\[4pt]
        \textbf{Custo}: $O(1)$ após encontrar raízes
      \end{block}
    \end{column}
    \begin{column}{0.47\textwidth}
      \begin{tikzpicture}
        %% Antes
        \node[ufroot] (ra) at (0,2.8) {1};
        \node[ufnode] (na2) at (-0.6,1.8) {2};
        \node[ufnode] (na3) at (0.6,1.8) {3};
        \draw[ufedge] (na2)--(ra); \draw[ufedge] (na3)--(ra);
        \node[ufroot] (rb) at (2.8,2.8) {4};
        \node[ufnode] (nb5) at (2.8,1.8) {5};
        \draw[ufedge] (nb5)--(rb);
        \node[font=\tiny,UFSGray] at (1.4,3.2) {antes};
        %% Depois
        \draw[-{Stealth},thick,UFSGray] (1.4,1.2)--(1.4,0.6);
        \node[ufroot] (rc) at (0.8,0) {1};
        \node[ufnode] (nc2) at (0,- 0.9) {2};
        \node[ufnode] (nc3) at (0.8,-0.9) {3};
        \node[ufnode] (nc4) at (1.8,-0.9) {4};
        \node[ufnode] (nc5) at (2.4,-1.8) {5};
        \draw[ufedge] (nc2)--(rc); \draw[ufedge] (nc3)--(rc);
        \draw[ufedge] (nc4)--(rc); \draw[ufedge] (nc5)--(nc4);
        \node[font=\tiny,UFSGray] at (0.8,-2.4)
          {\textsc{Union}(1,4)};
      \end{tikzpicture}
    \end{column}
  \end{columns}
\end{frame}

\begin{frame}{Sequência completa de operações}%% 22
  \begin{center}
  \begin{tikzpicture}[scale=0.92]
    %% Estado 0: singletons 1..6
    \node[font=\scriptsize\bfseries,UFSBlue,anchor=east] at (-0.35,4.2) {Início:};
    \foreach \i in {1,...,6}
      \node[ufroot,minimum size=0.55cm,font=\scriptsize] at ({(\i-1)*1.0},4.2) {\i};

    %% Estado 1: Union(1,2), Union(3,4), Union(5,6)
    \node[font=\scriptsize\bfseries,UFSBlue] at (-0.5,2.6) {U(1,2)};
    \node[font=\scriptsize\bfseries,UFSBlue] at (-0.5,2.2) {U(3,4)};
    \node[font=\scriptsize\bfseries,UFSBlue] at (-0.5,1.8) {U(5,6):};

    \node[ufroot,minimum size=0.52cm,font=\scriptsize] (s1r) at (0.5,2.6) {1};
    \node[ufnode,minimum size=0.52cm,font=\scriptsize] (s1a) at (0.5,1.8) {2};
    \draw[ufedge] (s1a)--(s1r);

    \node[ufroot,minimum size=0.52cm,font=\scriptsize] (s3r) at (2.2,2.6) {3};
    \node[ufnode,minimum size=0.52cm,font=\scriptsize] (s3a) at (2.2,1.8) {4};
    \draw[ufedge] (s3a)--(s3r);

    \node[ufroot,minimum size=0.52cm,font=\scriptsize] (s5r) at (3.9,2.6) {5};
    \node[ufnode,minimum size=0.52cm,font=\scriptsize] (s5a) at (3.9,1.8) {6};
    \draw[ufedge] (s5a)--(s5r);

    %% Estado 2: Union(1,3) — label recuado para não sobrepor nó 2
    \node[font=\scriptsize\bfseries,UFSBlue] at (-0.8,0.5) {U(1,3):};
    \node[ufroot,minimum size=0.52cm,font=\scriptsize] (t1r) at (0.5,0.8) {1};
    \node[ufnode,minimum size=0.52cm,font=\scriptsize] (t1a) at (0.0,0) {2};
    \node[ufnode,minimum size=0.52cm,font=\scriptsize] (t1b) at (0.9,0) {3};
    \node[ufnode,minimum size=0.52cm,font=\scriptsize] (t1c) at (1.5,-0.7) {4};
    \draw[ufedge] (t1a)--(t1r); \draw[ufedge] (t1b)--(t1r);
    \draw[ufedge] (t1c)--(t1b);

    %% Estado 3: Union(1,5)
    \node[font=\scriptsize\bfseries,UFSBlue] at (4.5,0.5) {U(1,5):};
    \node[ufroot,minimum size=0.52cm,font=\scriptsize] (u1r) at (6,1.0) {1};
    \node[ufnode,minimum size=0.52cm,font=\scriptsize] (u1a) at (5.2,0.1) {2};
    \node[ufnode,minimum size=0.52cm,font=\scriptsize] (u1b) at (5.8,0.1) {3};
    \node[ufnode,minimum size=0.52cm,font=\scriptsize] (u1c) at (6.6,0.1) {5};
    \node[ufnode,minimum size=0.52cm,font=\scriptsize] (u1d) at (5.8,-0.7) {4};
    \node[ufnode,minimum size=0.52cm,font=\scriptsize] (u1e) at (7.2,-0.7) {6};
    \draw[ufedge] (u1a)--(u1r); \draw[ufedge] (u1b)--(u1r);
    \draw[ufedge] (u1c)--(u1r);
    \draw[ufedge] (u1d)--(u1b); \draw[ufedge] (u1e)--(u1c);

    %% Setas entre estados
    \draw[-{Stealth},thick,UFSGray] (2,3.8)--(2,3.2);
    \draw[-{Stealth},thick,UFSGray] (2,1.4)--(2,0.8) node[right,font=\tiny] {};
    \draw[-{Stealth},thick,UFSGray] (3.2,0.5)--(4.4,0.5);
  \end{tikzpicture}
  \end{center}
\end{frame}

\begin{frame}{Problema: degeneração em cadeia}%% 23
  \begin{alertblock}{Sequência adversarial para a floresta ingênua}
    \textsc{Union}$(x_1,x_2)$, \textsc{Union}$(x_2,x_3)$, \ldots, \textsc{Union}$(x_{n-1},x_n)$
    sempre conectando a raiz da árvore maior como filho da menor.
  \end{alertblock}
  \vspace{0.3cm}
  \begin{center}
  \begin{tikzpicture}
    \node[ufroot] (n1) at (0,0) {$x_1$};
    \node[ufnode] (n2) at (1.2,0) {$x_2$};
    \node[ufnode] (n3) at (2.4,0) {$x_3$};
    \node at (3.6,0) {\ldots};
    \node[ufnode] (nn) at (4.8,0) {$x_n$};
    \draw[ufedge] (n2)--(n1); \draw[ufedge] (n3)--(n2);
    \draw[ufedge] (n3)--(3.2,0); \draw[ufedge] (4.0,0)--(nn);
    \node[font=\small,UFSGray,below=0.4cm of n3]
      {altura $= n-1$ → \textsc{Find-Set} em $O(n)$};
  \end{tikzpicture}
  \end{center}
\end{frame}

\begin{frame}{Comparação: lista vs.\ floresta (sem heurísticas)}%% 24
  \begin{center}
  \renewcommand{\arraystretch}{1.5}
  \begin{tabular}{lcc}
    \toprule
    & \textbf{Lista encadeada} & \textbf{Floresta ingênua} \\
    \midrule
    \textsc{Make-Set} & $O(1)$ & $O(1)$ \\
    \textsc{Find-Set} & $O(1)$ & $O(n)$ \\
    \textsc{Union}    & $O(n)$ & $O(1)$\,+\,\textsc{Find-Set} \\
    Total $m$ ops    & $O(mn)$ pior caso & $O(mn)$ pior caso \\
    \bottomrule
  \end{tabular}
  \end{center}
  \vspace{0.3cm}
  \begin{center}
    Ambas são equivalentes no pior caso. \\
    As \textbf{heurísticas} fazem a diferença — próximas seções.
  \end{center}
\end{frame}

\begin{frame}[fragile,shrink=15]{Implementação Java: floresta básica}%% 25
\begin{lstlisting}
public class UnionFind {
    private final int[] parent;

    public UnionFind(int n) {
        parent = new int[n];
        for (int i = 0; i < n; i++) parent[i] = i; // Make-Set
    }

    // Find-Set sem compressao -- O(altura)
    public int find(int x) {
        while (parent[x] != x) x = parent[x];
        return x;
    }

    // Union sem heuristica -- O(1) apos find
    public void union(int x, int y) {
        int rx = find(x), ry = find(y);
        if (rx != ry) parent[ry] = rx;
    }
}
\end{lstlisting}
\end{frame}

\begin{frame}{Recapitulando: Floresta}%% 26
  \begin{block}{O que aprendemos}
    \begin{itemize}
      \item Floresta: cada nó guarda apenas ponteiro para o pai
      \item \textsc{Make-Set} e \textsc{Union} em $O(1)$; \textsc{Find-Set} em $O(\text{altura})$
      \item Sem heurísticas, árvore degenera em cadeia → $O(n)$ por Find
      \item Solução: \textbf{Union by Rank} controla a altura da árvore
    \end{itemize}
  \end{block}
\end{frame}

%% ================================================================
%%  SEÇÃO 4 — UNION BY RANK
%% ================================================================
\section{Union by Rank}

\begin{frame}{Motivação: controlar a altura}%% 27
  \begin{columns}[T]
    \begin{column}{0.52\textwidth}
      \begin{block}{Problema}
        Sem critério de escolha, Union pode criar árvores altas.\\
        Find-Set percorre todo o caminho até a raiz.
      \end{block}
      \vspace{0.3cm}
      \begin{block}{Solução: Union by Rank}
        Manter um \textbf{rank} por nó — limite superior da altura.\\
        No Union, a raiz de menor rank passa a ser filha da de maior rank.
      \end{block}
    \end{column}
    \begin{column}{0.45\textwidth}
      \begin{tikzpicture}
        \node[font=\scriptsize,UFSRed] at (0.8,3.2) {ruim (degenerado)};
        \node[ufroot] (d1) at (0,2.6) {1}; \node[ufnode] (d2) at (1,2.6) {2};
        \node[ufnode] (d3) at (2,2.6) {3}; \draw[ufedge] (d2)--(d1); \draw[ufedge] (d3)--(d2);

        \node[font=\scriptsize,UFSBlue] at (0.8,1.6) {bom (balanceado)};
        \node[ufroot] (b1) at (0.8,1.0) {1};
        \node[ufnode] (b2) at (0,0.2) {2}; \node[ufnode] (b3) at (1.6,0.2) {3};
        \draw[ufedge] (b2)--(b1); \draw[ufedge] (b3)--(b1);
      \end{tikzpicture}
    \end{column}
  \end{columns}
\end{frame}

\begin{frame}{Definição formal de rank (CLRS, Seção 19.3)}%% 28
  \begin{columns}[T]
    \begin{column}{0.52\textwidth}
      \begin{block}{Rank}
        \begin{itemize}
          \item \textsc{Make-Set}$(x)$: $x.\texttt{rank} = 0$
          \item Se $\texttt{rank}[r_x] > \texttt{rank}[r_y]$: $r_y \to r_x$ (rank inalterado)
          \item Se ranks iguais: $r_y \to r_x$ e $r_x.\texttt{rank} \mathrel{+}= 1$
        \end{itemize}
      \end{block}
      \begin{alertblock}{Invariante}
        rank$(x)$ $\geq$ altura real da subárvore de $x$.\\
        Árvore de rank $k$ $\Rightarrow$ $\geq 2^k$ nós.
      \end{alertblock}
    \end{column}
    \begin{column}{0.45\textwidth}
      \begin{tikzpicture}[
        nr/.style={circle,draw=UFSBlue,fill=UFSBlue!15,minimum size=0.58cm,font=\scriptsize\bfseries},
        rt/.style={circle,draw=UFSRed,fill=UFSRed!20,very thick,minimum size=0.58cm,font=\scriptsize\bfseries},
        rl/.style={font=\tiny,UFSRed}
      ]
        \node[rt,label={[rl]right:r=3}] (r) at (1.8,3.2) {1};
        \node[nr,label={[rl]right:r=2}] (a) at (0.8,2.2) {2};
        \node[nr,label={[rl]right:r=2}] (b) at (2.8,2.2) {5};
        \node[nr,label={[rl]left:r=1}]  (c) at (0.0,1.2) {3};
        \node[nr,label={[rl]right:r=1}] (d) at (1.6,1.2) {4};
        \node[nr,label={[rl]right:r=1}] (e) at (2.8,1.2) {6};
        \node[nr,label={[rl]left:r=0}]  (f) at (0.0,0.2) {7};
        \node[nr,label={[rl]right:r=0}] (g) at (1.6,0.2) {8};
        \draw[ufedge] (a)--(r); \draw[ufedge] (b)--(r);
        \draw[ufedge] (c)--(a); \draw[ufedge] (d)--(a);
        \draw[ufedge] (e)--(b);
        \draw[ufedge] (f)--(c); \draw[ufedge] (g)--(d);
        \node[font=\tiny,UFSGray] at (1.8,-0.3) {altura real = 3 = rank da raiz};
      \end{tikzpicture}
    \end{column}
  \end{columns}
\end{frame}

\begin{frame}{Algoritmo Union by Rank}%% 29
  \begin{columns}[T]
    \begin{column}{0.55\textwidth}
      \begin{block}{\textsc{Link}$(r_x, r_y)$}
        \textbf{if} $r_x.\texttt{rank} > r_y.\texttt{rank}$:\\
        \quad $r_y.\texttt{parent} \leftarrow r_x$\\
        \textbf{elseif} $r_x.\texttt{rank} < r_y.\texttt{rank}$:\\
        \quad $r_x.\texttt{parent} \leftarrow r_y$\\
        \textbf{else}:\\
        \quad $r_y.\texttt{parent} \leftarrow r_x$\\
        \quad $r_x.\texttt{rank} \mathrel{+}= 1$
      \end{block}
      \begin{block}{\textsc{Union}$(x,y)$}
        \textsc{Link}(\textsc{Find-Set}$(x)$, \textsc{Find-Set}$(y)$)
      \end{block}
    \end{column}
    \begin{column}{0.42\textwidth}
      \begin{exampleblock}{Resultado}
        Altura máxima da árvore: $\lfloor\lg n\rfloor$\\[4pt]
        \textsc{Find-Set} garante $O(\lg n)$ com esta heurística isolada.
      \end{exampleblock}
    \end{column}
  \end{columns}
\end{frame}

\begin{frame}{Passo a passo: Union by Rank — Parte 1}%% 30
  \begin{center}
  \begin{tikzpicture}[
    nr/.style={ufnode,font=\scriptsize\bfseries,minimum size=0.55cm},
    rr/.style={ufroot,font=\scriptsize\bfseries,minimum size=0.55cm},
    lbl/.style={font=\scriptsize,UFSGray}
  ]
    \node[lbl] at (-0.5,3.5) {Início (r=0):};
    \foreach \i in {1,...,8}
      \node[rr,label={[font=\tiny]below:r=0}] at ({(\i-1)*1.0},3.5) {\i};

    \node[lbl] at (-0.5,2.0) {U(1,2):};
    \node[rr,label={[font=\tiny]below:r=1}] (u12r) at (0.5,2.2) {1};
    \node[nr] (u12a) at (0.5,1.4) {2};
    \draw[ufedge] (u12a)--(u12r);

    \node[lbl] at (2.0,2.0) {U(3,4):};
    \node[rr,label={[font=\tiny]below:r=1}] (u34r) at (3.0,2.2) {3};
    \node[nr] (u34a) at (3.0,1.4) {4};
    \draw[ufedge] (u34a)--(u34r);

    \node[lbl] at (4.5,2.0) {U(5,6):};
    \node[rr,label={[font=\tiny]below:r=1}] (u56r) at (5.5,2.2) {5};
    \node[nr] (u56a) at (5.5,1.4) {6};
    \draw[ufedge] (u56a)--(u56r);

    \node[lbl] at (7.0,2.0) {U(7,8):};
    \node[rr,label={[font=\tiny]below:r=1}] (u78r) at (8.0,2.2) {7};
    \node[nr] (u78a) at (8.0,1.4) {8};
    \draw[ufedge] (u78a)--(u78r);
  \end{tikzpicture}
  \end{center}
\end{frame}

\begin{frame}[shrink=20]{Passo a passo: Union by Rank — Parte 2}%% 31
  \begin{center}
  \begin{tikzpicture}[
    nr/.style={ufnode,font=\scriptsize\bfseries,minimum size=0.55cm},
    rr/.style={ufroot,font=\scriptsize\bfseries,minimum size=0.55cm},
    lbl/.style={font=\scriptsize,UFSGray}
  ]
    \node[lbl] at (-0.5,2.5) {U(1,3) r=1=r=1:};
    \node[rr,label={[font=\tiny]above:r=2}] (r1) at (1.0,2.5) {1};
    \node[nr] (a2) at (0.2,1.5) {2};
    \node[nr] (a3) at (1.0,1.5) {3};
    \node[nr] (a4) at (1.8,0.7) {4};
    \draw[ufedge] (a2)--(r1); \draw[ufedge] (a3)--(r1);
    \draw[ufedge] (a4)--(a3);

    \node[lbl] at (3.5,2.5) {U(5,7) r=1=r=1:};
    \node[rr,label={[font=\tiny]above:r=2}] (r5) at (5.0,2.5) {5};
    \node[nr] (b6) at (4.2,1.5) {6};
    \node[nr] (b7) at (5.0,1.5) {7};
    \node[nr] (b8) at (5.8,0.7) {8};
    \draw[ufedge] (b6)--(r5); \draw[ufedge] (b7)--(r5);
    \draw[ufedge] (b8)--(b7);

    \node[lbl] at (-0.5,-0.5) {U(1,5) r=2=r=2:};
    \node[rr,label={[font=\tiny]above:r=3}] (rr) at (4.0,-0.5) {1};
    \node[nr] (c2) at (3.0,-1.5) {2};
    \node[nr] (c3) at (3.8,-1.5) {3};
    \node[nr] (c5) at (4.6,-1.5) {5};
    \node[nr] (c4) at (3.8,-2.4) {4};
    \node[nr] (c6) at (4.6,-2.4) {6};
    \node[nr] (c7) at (5.4,-2.4) {7};
    \node[nr] (c8) at (5.4,-3.2) {8};
    \draw[ufedge] (c2)--(rr); \draw[ufedge] (c3)--(rr); \draw[ufedge] (c5)--(rr);
    \draw[ufedge] (c4)--(c3); \draw[ufedge] (c6)--(c5); \draw[ufedge] (c7)--(c5);
    \draw[ufedge] (c8)--(c7);
    \node[lbl] at (4.0,-3.8) {rank cresce apenas quando ranks são iguais};
  \end{tikzpicture}
  \end{center}
\end{frame}

\begin{frame}{Propriedades do rank}%% 32
  \begin{block}{Lema 19.2 (CLRS)}
    Para qualquer nó $x$: a subárvore de $x$ tem pelo menos $2^{x.\texttt{rank}}$ nós.
  \end{block}
  \vspace{0.3cm}
  \begin{block}{Corolário}
    O rank máximo de qualquer nó numa floresta com $n$ elementos é $\lfloor\lg n\rfloor$.\\
    Logo, a altura de qualquer árvore é no máximo $\lfloor\lg n\rfloor$.
  \end{block}
  \vspace{0.3cm}
  \begin{exampleblock}{Consequência}
    Com Union by Rank isolado: \textsc{Find-Set} em $O(\lg n)$.\\
    Para $n = 10^9$: altura $\leq 30$ — razoável, mas ainda podemos fazer melhor.
  \end{exampleblock}
\end{frame}

\begin{frame}{Prova do lema (esboço)}%% 33
  \begin{block}{Por que árvore de rank $k$ tem $\geq 2^k$ nós?}
    \textbf{Base} ($k=0$): rank 0 → subárvore tem $\geq 1 = 2^0$ nó. \checkmark\\[4pt]
    \textbf{Indução} ($k>0$): o rank de $r$ chega a $k$ somente quando \textsc{Link} une duas raízes de rank $k-1$.\\
    Por hipótese, cada uma tem $\geq 2^{k-1}$ nós.\\
    A árvore resultante tem $\geq 2^{k-1} + 2^{k-1} = 2^k$ nós. \checkmark
  \end{block}
  \vspace{0.2cm}
  \begin{alertblock}{Propriedade adicional}
    O número de nós com rank exatamente $k$ é $\leq n/2^k$.\\
    Os ranks são \textbf{estritamente crescentes} no caminho de qualquer nó até a raiz.
  \end{alertblock}
\end{frame}

\begin{frame}[fragile,shrink=25]{Código Java: Union by Rank}%% 34
\begin{lstlisting}
public class UnionFindRank {
    private final int[] parent;
    private final int[] rank;

    public UnionFindRank(int n) {
        parent = new int[n];
        rank   = new int[n];
        for (int i = 0; i < n; i++) parent[i] = i; // Make-Set
    }

    public int find(int x) {                 // Find-Set sem compressao
        while (parent[x] != x) x = parent[x];
        return x;
    }

    public void union(int x, int y) {        // Union by Rank
        int rx = find(x), ry = find(y);
        if (rx == ry) return;
        if (rank[rx] < rank[ry]) { int t = rx; rx = ry; ry = t; }
        parent[ry] = rx;                     // ry passa a ser filho de rx
        if (rank[rx] == rank[ry]) rank[rx]++;
    }
}
\end{lstlisting}
\end{frame}

\begin{frame}{Recapitulando: Union by Rank}%% 35
  \begin{block}{O que aprendemos}
    \begin{itemize}
      \item Rank: limite superior da altura do nó
      \item Union sempre coloca a raiz de menor rank como filha
      \item Árvore de rank $k$ tem pelo menos $2^k$ nós
      \item Altura máxima: $\lfloor\lg n\rfloor$ $\Rightarrow$ \textsc{Find-Set} em $O(\lg n)$
      \item Próximo passo: \textbf{Path Compression} para reduzir ainda mais
    \end{itemize}
  \end{block}
\end{frame}

%% ================================================================
%%  SEÇÃO 5 — PATH COMPRESSION
%% ================================================================
\section{Path Compression}

\begin{frame}[shrink=10]{O problema dos caminhos longos}%% 36
  \begin{columns}[T]
    \begin{column}{0.5\textwidth}
      \begin{alertblock}{Custo do Find-Set}
        Mesmo com Union by Rank, \textsc{Find-Set}$(x)$ percorre todo o caminho de $x$ até a raiz.\\[4pt]
        Chamadas repetidas ao mesmo nó percorrem o mesmo caminho inutilmente.
      \end{alertblock}
      \vspace{0.3cm}
      \begin{block}{Ideia}
        Durante \textsc{Find-Set}, \textbf{comprimir o caminho}: fazer todos os nós visitados apontarem diretamente para a raiz.
      \end{block}
    \end{column}
    \begin{column}{0.47\textwidth}
      \begin{tikzpicture}
        \node[ufroot] (r) at (1.2,4.0) {1};
        \node[ufnode] (a) at (1.2,3.0) {2};
        \node[ufnode] (b) at (1.2,2.0) {3};
        \node[ufnode] (c) at (1.2,1.0) {4};
        \node[ufnode] (d) at (1.2,0.0) {5};
        \draw[ufedge] (a)--(r); \draw[ufedge] (b)--(a);
        \draw[ufedge] (c)--(b); \draw[ufedge] (d)--(c);
        \node[font=\tiny,UFSRed,align=center] at (2.4,2.0)
          {Find-Set(5)\\percorre 4 arestas};
      \end{tikzpicture}
    \end{column}
  \end{columns}
\end{frame}

\begin{frame}{Algoritmo: Find-Set com Path Compression}%% 37
  \begin{columns}[T]
    \begin{column}{0.5\textwidth}
      \begin{block}{\textsc{Find-Set} recursivo (two-pass)}
        \textbf{if} $x.\texttt{parent} \neq x$:\\
        \quad $x.\texttt{parent} \leftarrow \textsc{Find-Set}(x.\texttt{parent})$\\
        \textbf{return} $x.\texttt{parent}$
      \end{block}
      \vspace{0.3cm}
      \begin{block}{\textsc{Find-Set} iterativo}
        Primeiro percurso: encontrar a raiz.\\
        Segundo percurso: atualizar todos os ponteiros para a raiz.
      \end{block}
    \end{column}
    \begin{column}{0.47\textwidth}
      \begin{exampleblock}{O que muda?}
        \begin{itemize}
          \item O valor retornado é idêntico ao sem compressão
          \item Os ponteiros \texttt{parent} são \textbf{atualizados como efeito colateral}
          \item O rank \textbf{não é atualizado} (pode ficar maior que a altura real)
          \item Custo amortizado: $O(\lg n)$ com PC isolado
        \end{itemize}
      \end{exampleblock}
    \end{column}
  \end{columns}
\end{frame}

\begin{frame}{Passo a passo: antes da compressão}%% 38
  \begin{center}
  \begin{tikzpicture}
    \node[font=\small\bfseries,UFSBlue] at (-1.5,2) {Antes:};
    \node[ufroot] (r)  at (0,4) {1};
    \node[ufnode] (n2) at (0,3) {2};
    \node[ufnode] (n3) at (0,2) {3};
    \node[ufnode] (n4) at (0,1) {4};
    \node[ufnode] (n5) at (0,0) {5};
    \draw[ufedge] (n2)--(r);
    \draw[ufedge] (n3)--(n2);
    \draw[ufedge] (n4)--(n3);
    \draw[ufedge] (n5)--(n4);

    %% Seta de operacao
    \draw[-{Stealth},thick,UFSGray] (1.5,2)--(2.5,2);
    \node[font=\small] at (2.0,2.5) {\textsc{Find-Set}(5)};

    %% Caminho visitado
    \draw[-{Stealth[scale=1.0]},very thick,UFSYellow!80!black,bend right=30]
      (n5) to (n4);
    \draw[-{Stealth[scale=1.0]},very thick,UFSYellow!80!black,bend right=30]
      (n4) to (n3);
    \draw[-{Stealth[scale=1.0]},very thick,UFSYellow!80!black,bend right=30]
      (n3) to (n2);
    \draw[-{Stealth[scale=1.0]},very thick,UFSYellow!80!black,bend right=30]
      (n2) to (r);
    \node[font=\tiny,UFSGray] at (-1.5,0) {percorre 4 arestas};
  \end{tikzpicture}
  \end{center}
\end{frame}

\begin{frame}{Passo a passo: após a compressão}%% 39
  \begin{center}
  \begin{tikzpicture}
    \node[font=\small\bfseries,UFSBlue] at (-2.5,1) {Após Find-Set(5):};
    \node[ufroot] (r) at (0,2) {1};
    \node[ufnode] (n2) at (-1.8,0.6) {2};
    \node[ufnode] (n3) at (-0.6,0.6) {3};
    \node[ufnode] (n4) at (0.6,0.6) {4};
    \node[ufnode] (n5) at (1.8,0.6) {5};
    \draw[ufedge,UFSBlue,very thick] (n2)--(r);
    \draw[ufedge,UFSBlue,very thick] (n3)--(r);
    \draw[ufedge,UFSBlue,very thick] (n4)--(r);
    \draw[ufedge,UFSBlue,very thick] (n5)--(r);
    \node[font=\small,UFSGray] at (0,-0.4)
      {todos os nós 2,3,4,5 agora apontam diretamente para a raiz 1};
    \node[font=\small,UFSBlue] at (0,-0.9)
      {próximo Find-Set: $O(1)$ por nó};
  \end{tikzpicture}
  \end{center}
\end{frame}

\begin{frame}[fragile,shrink=30]{Implementação Java: Path Compression recursivo}%% 40
\begin{lstlisting}
public class UnionFindPC {
    private final int[] parent;
    private final int[] rank;

    public UnionFindPC(int n) {
        parent = new int[n];
        rank   = new int[n];
        for (int i = 0; i < n; i++) parent[i] = i;
    }

    // Find-Set com path compression (recursivo)
    public int find(int x) {
        if (parent[x] != x)
            parent[x] = find(parent[x]); // comprime no retorno da recursao
        return parent[x];
    }

    public void union(int x, int y) {
        int rx = find(x), ry = find(y);
        if (rx == ry) return;
        if (rank[rx] < rank[ry]) { int t = rx; rx = ry; ry = t; }
        parent[ry] = rx;
        if (rank[rx] == rank[ry]) rank[rx]++;
    }
}
\end{lstlisting}
\end{frame}

\begin{frame}[fragile,shrink=25]{Implementação Java: Path Compression iterativo}%% 41
\begin{lstlisting}
// Dois percursos: encontrar raiz, depois comprimir
public int find(int x) {
    // 1. encontrar a raiz
    int root = x;
    while (parent[root] != root) root = parent[root];

    // 2. comprimir: todos apontam para root
    while (parent[x] != root) {
        int next = parent[x];
        parent[x] = root;
        x = next;
    }
    return root;
}
\end{lstlisting}
\vspace{0.3cm}
\begin{block}{Comparação}
  \begin{itemize}
    \item Recursivo: elegante, mas usa pilha de chamadas
    \item Iterativo: preferível para pilhas rasas e grandes $n$
    \item Ambos produzem o mesmo resultado final
  \end{itemize}
\end{block}
\end{frame}

\begin{frame}{Recapitulando: Path Compression}%% 42
  \begin{block}{O que aprendemos}
    \begin{itemize}
      \item PC faz todos os nós no caminho apontarem diretamente para a raiz
      \item O rank \emph{não} é corrigido — pode ser maior que a altura real
      \item Custo amortizado com PC isolado: $O(\lg n)$ por operação Find
      \item PC e Union by Rank \emph{juntos} → $O(\alpha(n))$ — próxima seção
    \end{itemize}
  \end{block}
\end{frame}

%% ================================================================
%%  SEÇÃO 6 — COMBINAÇÃO DAS HEURÍSTICAS
%% ================================================================
\section{Combinação: Rank + Path Compression}

\begin{frame}{Por que as heurísticas se complementam?}%% 43
  \begin{columns}[T]
    \begin{column}{0.5\textwidth}
      \begin{block}{Union by Rank}
        Controla a \textbf{estrutura de longo prazo}:\\
        limita a altura durante os Unions.
      \end{block}
      \vspace{0.3cm}
      \begin{block}{Path Compression}
        Reduz o custo de \textbf{futuras consultas}:\\
        achata a árvore durante os Finds.
      \end{block}
    \end{column}
    \begin{column}{0.47\textwidth}
      \begin{exampleblock}{Sinergia}
        PC achata nós visitados → as futuras Finds são baratas.\\[6pt]
        Rank preserva a propriedade de que as raízes têm rank alto → PC só ocorre em caminhos curtos (amortizado).\\[6pt]
        Resultado: $O(\alpha(n))$ amortizado por operação.
      \end{exampleblock}
    \end{column}
  \end{columns}
\end{frame}

\begin{frame}{Passo a passo combinado — Parte 1}%% 44
  \begin{center}
  \begin{tikzpicture}[
    nr/.style={ufnode,font=\scriptsize\bfseries,minimum size=0.55cm},
    rr/.style={ufroot,font=\scriptsize\bfseries,minimum size=0.55cm},
    lbl/.style={font=\scriptsize,UFSGray}
  ]
    \node[lbl] at (-0.5,3.2) {Após U(1,2), U(3,4), U(1,3):};
    \node[rr,label={[font=\tiny]above:r=2}] (rt) at (2,3.2) {1};
    \node[nr] (n2) at (1,2.2) {2};
    \node[nr] (n3) at (2,2.2) {3};
    \node[nr] (n4) at (3,1.4) {4};
    \draw[ufedge] (n2)--(rt); \draw[ufedge] (n3)--(rt);
    \draw[ufedge] (n4)--(n3);

    \draw[-{Stealth},thick,UFSGray] (4.5,2.6)--(5.2,2.6);
    \node[lbl] at (6.8,3.2) {Find-Set(4):};
    \node[rr,label={[font=\tiny]above:r=2}] (rt2) at (7.5,3.2) {1};
    \node[nr] (n2b) at (6.5,2.2) {2};
    \node[nr] (n3b) at (7.5,2.2) {3};
    \node[nr] (n4b) at (8.5,2.2) {4};
    \draw[ufedge] (n2b)--(rt2); \draw[ufedge] (n3b)--(rt2);
    \draw[ufedge,UFSBlue,very thick] (n4b)--(rt2);
    \node[lbl] at (7.5,1.4) {4 comprimido: aponta para 1};
  \end{tikzpicture}
  \end{center}
\end{frame}

\begin{frame}{Passo a passo combinado — Parte 2}%% 45
  \begin{center}
  \begin{tikzpicture}[
    nr/.style={ufnode,font=\scriptsize\bfseries,minimum size=0.55cm},
    rr/.style={ufroot,font=\scriptsize\bfseries,minimum size=0.55cm},
    lbl/.style={font=\scriptsize,UFSGray}
  ]
    \node[lbl] at (0,3.8) {Adicionar U(5,6), U(7,8), U(5,7), U(1,5):};
    \node[rr,label={[font=\tiny]above:r=3}] (rt) at (4,3.2) {1};
    \node[nr] (n2) at (1.8,2.2) {2};
    \node[nr] (n3) at (3,2.2) {3};
    \node[nr] (n5) at (4.5,2.2) {5};
    \node[nr] (n4) at (3,1.2) {4};
    \node[nr] (n6) at (3.8,1.2) {6};
    \node[nr] (n7) at (5.2,1.2) {7};
    \node[nr] (n8) at (6,0.4) {8};
    \draw[ufedge] (n2)--(rt); \draw[ufedge] (n3)--(rt);
    \draw[ufedge] (n5)--(rt);
    \draw[ufedge] (n4)--(n3); \draw[ufedge] (n6)--(n5);
    \draw[ufedge] (n7)--(n5); \draw[ufedge] (n8)--(n7);

    \draw[-{Stealth},thick,UFSGray] (7.0,2.2)--(7.8,2.2);
    \node[lbl] at (9,3.4) {Find-Set(8):};
    \node[rr,label={[font=\tiny]above:r=3}] (rt2) at (9,2.6) {1};
    \node[nr] (n5b) at (8.2,1.6) {5};
    \node[nr] (n7b) at (9.0,1.6) {7};
    \node[nr] (n8b) at (9.8,1.6) {8};
    \draw[ufedge] (n5b)--(rt2);
    \draw[ufedge,UFSBlue,very thick] (n7b)--(rt2);
    \draw[ufedge,UFSBlue,very thick] (n8b)--(rt2);
    \node[lbl] at (9,0.8) {7 e 8 comprimidos};
  \end{tikzpicture}
  \end{center}
\end{frame}

\begin{frame}[fragile,shrink=5]{Código Java: GenericUnionFind\textless T\textgreater{} — Parte 1}%% 46
\begin{lstlisting}
import java.util.HashMap;
import java.util.Map;

public class GenericUnionFind<T> {
    private final Map<T, T>       parent = new HashMap<>();
    private final Map<T, Integer> rank   = new HashMap<>();

    public void makeSet(T x) {
        parent.put(x, x);
        rank.put(x, 0);
    }

    public T find(T x) {        // pre-condicao: makeSet(x) ja chamado
        if (!parent.get(x).equals(x))
            parent.put(x, find(parent.get(x))); // path compression
        return parent.get(x);
    }
\end{lstlisting}
\end{frame}

\begin{frame}[fragile,shrink=5]{Código Java: GenericUnionFind\textless T\textgreater{} — Parte 2}%% 47
\begin{lstlisting}
    public boolean union(T x, T y) {
        T rx = find(x), ry = find(y);
        if (rx.equals(ry)) return false; // ja no mesmo conjunto

        int rankX = rank.get(rx), rankY = rank.get(ry);
        if (rankX < rankY) { T tmp = rx; rx = ry; ry = tmp; }

        parent.put(ry, rx);              // ry passa a ser filho de rx
        if (rankX == rankY)              // ranks iguais: incrementar
            rank.merge(rx, 1, Integer::sum);
        return true;
    }

    public boolean connected(T x, T y) {
        return find(x).equals(find(y));
    }
}
\end{lstlisting}
\end{frame}

\begin{frame}{Recapitulando: Combinação}%% 48
  \begin{block}{O que aprendemos}
    \begin{itemize}
      \item Union by Rank controla a estrutura de longo prazo (altura $\leq \lg n$)
      \item Path Compression achata o caminho durante cada Find
      \item Juntos: custo amortizado $O(\alpha(n))$ por operação
      \item GenericUnionFind<T> implementa em Java com HashMap e generics
    \end{itemize}
  \end{block}
\end{frame}

%% ================================================================
%%  SEÇÃO 7 — ANÁLISE AMORTIZADA
%% ================================================================
\section{Análise Amortizada}

\begin{frame}{Por que análise amortizada?}%% 49
  \begin{block}{Custo por operação vs.\ custo total}
    \begin{itemize}
      \item Uma única \textsc{Find-Set} com PC pode ser cara ($O(\lg n)$)
      \item Mas operações subsequentes sobre os mesmos nós são baratas
      \item Análise por operação piora o limite real — análise \textbf{amortizada} é mais justa
    \end{itemize}
  \end{block}
  \vspace{0.3cm}
  \begin{exampleblock}{Análise amortizada pelo método do potencial}
    Atribuímos um \textbf{potencial} $\Phi$ ao estado da estrutura.\\
    Custo amortizado = custo real + $\Delta\Phi$.\\
    Somando sobre $m$ operações: custo total $\leq \sum \text{amortizado} + \Phi_0$.\\[4pt]
    \footnotesize A função $\Phi$ concreta (baseada em $\mathrm{level}$ e $\mathrm{iter}$) é definida adiante,
    após introduzirmos a função de Ackermann.
  \end{exampleblock}
\end{frame}

\begin{frame}{Função de Ackermann — intuição}%% 50
  \begin{block}{Uma função que cresce \emph{incrivelmente} rápido}
    \begin{itemize}
      \item $A_0(j) = j+1$ (sucessor);\quad $A_1(j) = 2j+1$
      \item $A_2(j) = 2^{j+1}(j+1)-1$;\quad logo $A_2(1) = 7$
      \item $A_3(1) = A_2(7) = 2047$
      \item $A_4(1) = A_3(2047)$ — maior que o nº de átomos no universo
    \end{itemize}
  \end{block}
  \vspace{0.3cm}
  \begin{alertblock}{Inversa de Ackermann: $\alpha(n)$}
    $\alpha(n) = \min\{k : A_k(1) \geq n\}$\\[4pt]
    Como $A_k(1)$ explode, $\alpha(n) \leq 4$ para \emph{qualquer} $n$ prático\\
    (basta $n \leq A_4(1)$, um número astronômico).
  \end{alertblock}
\end{frame}

\begin{frame}[shrink=20]{Função de Ackermann $A_k(j)$ (CLRS, Seção 19.4)}%% 51
  \begin{block}{Definição (nível $k\geq 0$, argumento $j\geq 1$)}
    \[
      A_k(j) =
      \begin{cases}
        j+1                 & \text{se } k = 0 \\[2pt]
        A_{k-1}^{(j+1)}(j)  & \text{se } k \geq 1
      \end{cases}
    \]
    onde $A_{k-1}^{(j+1)}$ é $A_{k-1}$ aplicada $j+1$ vezes (iteração funcional).
  \end{block}
  \vspace{0.2cm}
  \begin{center}
  \renewcommand{\arraystretch}{1.3}
  \begin{tabular}{ccl}
    \toprule
    $k$ & $A_k(1)$ & fórmula geral \\
    \midrule
    0 & 2    & $A_0(j) = j+1$ \\
    1 & 3    & $A_1(j) = 2j+1$ \\
    2 & 7    & $A_2(j) = 2^{j+1}(j+1)-1$ \\
    3 & 2047 & $A_3(1) = A_2(7)$ \\
    4 & \multicolumn{2}{l}{$A_4(1) = A_3(2047) \gg 10^{80}$ (astronômico)} \\
    \bottomrule
  \end{tabular}
  \end{center}
\end{frame}

\begin{frame}[shrink=20]{Inversa de Ackermann: $\alpha(n)$}%% 52
  \begin{block}{Definição}
    $\alpha(n) = \min\{\,k : A_k(1) \geq n\,\}$ \quad — o menor nível cujo $A_k(1)$ já alcança $n$.
  \end{block}
  \vspace{0.3cm}
  \begin{center}
  \renewcommand{\arraystretch}{1.3}
  \begin{tabular}{ll}
    \toprule
    \textbf{Faixa de $n$} & $\alpha(n)$ \\
    \midrule
    $n \leq 2$                 & 0 \\
    $n = 3$                    & 1 \\
    $4 \leq n \leq 7$          & 2 \\
    $8 \leq n \leq 2047$       & 3 \\
    $2048 \leq n \leq A_4(1)$  & 4 \quad(todo $n$ concebível) \\
    \bottomrule
  \end{tabular}
  \end{center}
  \vspace{0.3cm}
  \begin{center}
    Como $A_4(1) \gg 10^{80}$, na prática $\alpha(n) \leq 4$ sempre.
  \end{center}
\end{frame}

\begin{frame}{Crescimento comparado}%% 53
  \begin{center}
  \begin{tikzpicture}[xscale=0.9,yscale=0.7]
    % Axes
    \draw[-{Stealth},thick] (0,0)--(8.5,0) node[right] {$n$};
    \draw[-{Stealth},thick] (0,0)--(0,5.5) node[above] {custo};
    % Grid
    \foreach \y in {1,...,5}
      \draw[dotted,UFSGray!40] (0,\y)--(8,\y);
    % O(n) linear
    \draw[thick,UFSRed]   (0,0)--(7,5)   node[right,font=\scriptsize] {$O(n)$};
    % O(n log n)
    \draw[thick,UFSYellow!80!black] (0,0) plot[domain=1:4,samples=30]
      (\x*1.75,{\x*log2(\x)/4}) node[right,font=\scriptsize] {$O(n\lg n)$};
    % O(n alpha(n))
    \draw[thick,UFSBlue,very thick] (0,0)--(7,0.3)
      node[right,font=\scriptsize] {$O(n\,\alpha(n))$};
    % Labels
    \node[font=\scriptsize,UFSGray] at (4,-0.4) {escala logarítmica};
  \end{tikzpicture}
  \end{center}
\end{frame}

\begin{frame}{Teorema 19.14 (CLRS 4ª ed.)}%% 54
  \begin{block}{Enunciado}
    Seja uma sequência de $m$ operações \textsc{Make-Set}, \textsc{Union} e \textsc{Find-Set}
    sobre $n$ elementos, usando \textbf{union por rank} e \textbf{path compression}.\\[6pt]
    O tempo total no pior caso é:
    \[
      O(m\,\alpha(n))
    \]
  \end{block}
  \vspace{0.3cm}
  \begin{exampleblock}{Significado prático}
    Como $\alpha(n) \leq 4$ para qualquer $n$ prático, o custo por operação é essencialmente $O(1)$ amortizado.
  \end{exampleblock}
\end{frame}

\begin{frame}{Blocos de rank (intuição)}%% 55
  \begin{columns}[T]
    \begin{column}{0.5\textwidth}
      \begin{block}{Definição}
        $B_k = \bigl\{A_{k-1}(1)+1,\ldots,A_k(1)\bigr\}$,\ \ $B_0=\{0,\ldots,A_0(1)\}$\\[4pt]
        \renewcommand{\arraystretch}{1.25}
        \begin{tabular}{cl}
          \toprule
          \textbf{Bloco} & \textbf{Ranks} \\
          \midrule
          $B_0$ & $\{0,1,2\}$ \\
          $B_1$ & $\{3\}$ \\
          $B_2$ & $\{4,\ldots,7\}$ \\
          $B_3$ & $\{8,\ldots,2047\}$ \\
          $B_4$ & $\{2048,\ldots,A_4(1)\}$ \\
          \bottomrule
        \end{tabular}
      \end{block}
    \end{column}
    \begin{column}{0.47\textwidth}
      \begin{exampleblock}{Por que isso importa?}
        Em qualquer árvore com $n$ nós práticos:\\
        \begin{itemize}
          \item ranks $\leq \lfloor\lg n\rfloor \leq 30$ (para $n=10^9$)
          \item logo só os blocos $B_0$--$B_3$ são usados
          \item \textbf{$\leq 4$ blocos} $\Rightarrow$ $O(\alpha(n))$ saltos por Find-Set
        \end{itemize}
        Esta é a origem do $\alpha(n)$.
      \end{exampleblock}
    \end{column}
  \end{columns}
\end{frame}

\begin{frame}[shrink=15]{Por que surge $\alpha(n)$? (intuição)}%% 56
  \begin{columns}[T]
    \begin{column}{0.46\textwidth}
      \begin{block}{Dois tipos de passo de $x$ à raiz}
        Subindo, os ranks \textbf{crescem}. Agrupe-os em \textbf{blocos}; cada
        passo é de um tipo:\\[4pt]
        \textcolor{UFSRed}{\textbf{Salta de bloco}}: pai em bloco superior.\\[3pt]
        \textcolor{UFSGray}{\textbf{Mesmo bloco}}: pai no mesmo bloco.
      \end{block}
      \begin{exampleblock}{Contagem por Find}
        \begin{itemize}
          \item \textcolor{UFSRed}{Saltos de bloco}: $\leq$ nº de blocos $= O(\alpha(n))$
          \item \textcolor{UFSGray}{Mesmo bloco}: pago pelo potencial $\Rightarrow$ amort.\ $O(1)$
        \end{itemize}
      \end{exampleblock}
    \end{column}
    \begin{column}{0.52\textwidth}
      \centering
      \begin{tikzpicture}[
        nd/.style={circle,draw=UFSGray,fill=white,thick,
          minimum size=0.55cm,font=\scriptsize\bfseries,inner sep=0pt},
        rt/.style={circle,draw=UFSBlue,fill=UFSBlue!15,very thick,
          minimum size=0.55cm,font=\scriptsize\bfseries,inner sep=0pt},
        same/.style={-{Stealth[scale=0.7]},thick,UFSGray!70},
        cross/.style={-{Stealth[scale=0.9]},very thick,UFSRed},
      ]
        %% Faixas de rank (fundo)
        \fill[UFSBlue!8]    (-1.7,-0.5) rectangle (1.9,1.35);
        \fill[UFSYellow!22] (-1.7,1.35) rectangle (1.9,3.15);
        \fill[UFSRed!8]     (-1.7,3.15) rectangle (1.9,5.35);
        \node[font=\scriptsize\bfseries,UFSBlue,anchor=west]           at (-1.65,0.45) {$B_1$};
        \node[font=\scriptsize\bfseries,UFSYellow!70!black,anchor=west] at (-1.65,2.25) {$B_2$};
        \node[font=\scriptsize\bfseries,UFSRed,anchor=west]            at (-1.65,4.25) {$B_3$};
        %% direção do rank
        \draw[-{Stealth[scale=0.7]},UFSGray!60] (1.55,-0.3)--(1.55,5.1);
        \node[font=\tiny,UFSGray,rotate=-90] at (1.82,2.4) {rank cresce};
        %% Nós do caminho (folha -> raiz), ranks crescentes
        \node[nd] (a) at (0.3,0.0)  {$x$};
        \node[nd] (b) at (0.3,0.9)  {};
        \node[nd] (c) at (0.3,2.1)  {};
        \node[nd] (d) at (0.3,2.7)  {};
        \node[nd] (e) at (0.3,3.9)  {};
        \node[rt] (r) at (0.3,4.8)  {raiz};
        %% Arestas: mesmo bloco (cinza) vs salto de bloco (vermelho)
        \draw[same]  (a)--(b);
        \draw[cross] (b)--(c);
        \draw[same]  (c)--(d);
        \draw[cross] (d)--(e);
        \draw[same]  (e)--(r);
      \end{tikzpicture}\\[2pt]
      {\footnotesize\color{UFSRed}\textbf{2 saltos de bloco} $\Rightarrow O(\alpha(n))$ por Find}
    \end{column}
  \end{columns}
\end{frame}

\begin{frame}[shrink=28]{Exemplo numérico: blocos num caminho de Find}%% 56a
  \begin{columns}[T]
    \begin{column}{0.5\textwidth}
      \begin{block}{Caminho de $x$ até a raiz}
        Ranks ao subir (sempre crescentes):
        \[ 0 \to 1 \to 3 \to 6 \to 10 \to 14 \]
        \renewcommand{\arraystretch}{1.2}
        \begin{tabular}{cl}
          \toprule
          \textbf{Rank} & \textbf{Bloco} \\
          \midrule
          $0,\ 1$   & $B_0=\{0,1,2\}$ \\
          $3$       & $B_1=\{3\}$ \\
          $6$       & $B_2=\{4,\ldots,7\}$ \\
          $10,\ 14$ & $B_3=\{8,\ldots,2047\}$ \\
          \bottomrule
        \end{tabular}
      \end{block}
      \begin{exampleblock}{Contagem}
        \textcolor{UFSRed}{\textbf{3 saltos}} $+$
        \textcolor{UFSGray}{\textbf{2 no mesmo bloco}}.\\[2pt]
        Caros $= 3 = (\text{nº blocos}-1) \leq \alpha(n)$,
        qualquer que seja o tamanho do caminho.
      \end{exampleblock}
    \end{column}
    \begin{column}{0.48\textwidth}
      \centering
      \begin{tikzpicture}[
        nd/.style={circle,draw=UFSGray,fill=white,thick,
          minimum size=0.6cm,font=\scriptsize\bfseries,inner sep=0pt},
        rt/.style={circle,draw=UFSBlue,fill=UFSBlue!15,very thick,
          minimum size=0.6cm,font=\scriptsize\bfseries,inner sep=0pt},
        same/.style={-{Stealth[scale=0.7]},thick,UFSGray!70},
        cross/.style={-{Stealth[scale=0.9]},very thick,UFSRed},
      ]
        %% Faixas de rank (fundo)
        \fill[UFSBlue!8]    (-1.5,-0.45) rectangle (1.7,1.0);
        \fill[UFSGray!14]   (-1.5,1.0)   rectangle (1.7,2.05);
        \fill[UFSYellow!22] (-1.5,2.05)  rectangle (1.7,3.35);
        \fill[UFSRed!8]     (-1.5,3.35)  rectangle (1.7,5.4);
        \node[font=\scriptsize\bfseries,UFSBlue,anchor=west]            at (-1.45,0.3) {$B_0$};
        \node[font=\scriptsize\bfseries,UFSGray,anchor=west]            at (-1.45,1.5) {$B_1$};
        \node[font=\scriptsize\bfseries,UFSYellow!70!black,anchor=west] at (-1.45,2.7) {$B_2$};
        \node[font=\scriptsize\bfseries,UFSRed,anchor=west]             at (-1.45,4.3) {$B_3$};
        %% Nós do caminho com o rank dentro
        \node[nd] (a) at (0.4,0.0)  {0};
        \node[nd] (b) at (0.4,0.6)  {1};
        \node[nd] (c) at (0.4,1.5)  {3};
        \node[nd] (d) at (0.4,2.7)  {6};
        \node[nd] (e) at (0.4,4.1)  {10};
        \node[rt] (r) at (0.4,4.95) {14};
        \node[font=\tiny,UFSBlue,right=2pt of r] {raiz};
        %% Arestas
        \draw[same]  (a)--(b);
        \draw[cross] (b)--(c);
        \draw[cross] (c)--(d);
        \draw[cross] (d)--(e);
        \draw[same]  (e)--(r);
      \end{tikzpicture}\\[2pt]
      {\footnotesize\color{UFSRed}\textbf{3 saltos} (vermelhos) cobrem 4 blocos}
    \end{column}
  \end{columns}
\end{frame}

\begin{frame}[shrink=15]{A função potencial $\Phi$ (CLRS, Seção 19.4)}%% 56b
  \begin{block}{Potencial de um nó $x$ (após $q$ operações)}
    Para um nó $x$ que \emph{não} é raiz e tem $x.\texttt{rank}\geq 1$, definimos
    \[
      \mathrm{level}(x) = \max\{\,k : x.p.\texttt{rank} \geq A_k(x.\texttt{rank})\,\},
    \]
    \[
      \mathrm{iter}(x) = \max\{\,i : x.p.\texttt{rank} \geq A_{\mathrm{level}(x)}^{(i)}(x.\texttt{rank})\,\}.
    \]
    \vspace{-0.25cm}
    \[
      \phi_q(x) =
      \begin{cases}
        \alpha(n)\cdot x.\texttt{rank}
          & \text{se } x \text{ é raiz ou } x.\texttt{rank}=0,\\[4pt]
        \bigl(\alpha(n)-\mathrm{level}(x)\bigr)\,x.\texttt{rank} - \mathrm{iter}(x)
          & \text{caso contrário.}
      \end{cases}
    \]
  \end{block}
  \begin{exampleblock}{Potencial total da floresta}
    $\displaystyle \Phi_q = \sum_{x} \phi_q(x)$.\quad
    No início todos os ranks são $0$, logo $\Phi_0 = 0$.
  \end{exampleblock}
\end{frame}

\begin{frame}[shrink=12]{Por que esse potencial funciona}%% 56c
  \begin{columns}[T]
    \begin{column}{0.5\textwidth}
      \begin{block}{Propriedades (Lemas 19.11--19.12)}
        \begin{itemize}
          \item $0 \leq \mathrm{level}(x) < \alpha(n)$
          \item $1 \leq \mathrm{iter}(x) \leq x.\texttt{rank}$
          \item $0 \leq \phi_q(x) \leq \alpha(n)\cdot x.\texttt{rank}$
          \item Logo $\Phi_q \geq 0$ sempre, e $\Phi_0 = 0$
        \end{itemize}
      \end{block}
    \end{column}
    \begin{column}{0.5\textwidth}
      \begin{block}{Custo amortizado por operação}
        \begin{description}
          \item[\textsc{Make-Set}:] $\Delta\Phi = 0 \Rightarrow O(1)$
          \item[\textsc{Union}:] $O(\alpha(n))$
          \item[\textsc{Find-Set}:] $O(\alpha(n))$
        \end{description}
      \end{block}
    \end{column}
  \end{columns}
  \vspace{0.2cm}
  \begin{alertblock}{Ideia central da prova de Find-Set (Lema 19.13)}
    Na compressão, cada nó que \emph{não} é raiz, \emph{não} é filho da raiz e cujo
    $\mathrm{level}$ coincide com o do pai tem seu $\phi$ \textbf{estritamente decrescido}
    (aumenta $\mathrm{iter}$ ou $\mathrm{level}$). Sobram $\leq \alpha(n)+2$ nós que pagam custo
    real; os demais saem da queda de $\Phi$ $\Rightarrow$ amortizado $O(\alpha(n))$.
  \end{alertblock}
\end{frame}

\begin{frame}{Limite da análise amortizada}%% 57
  \begin{block}{Custo total de $m$ operações}
    \begin{description}
      \item[\textsc{Make-Set}:] $n$ operações, $O(1)$ cada → total $O(n)$
      \item[\textsc{Union}:] $\leq n-1$ operações, $O(1)$ cada → total $O(n)$
      \item[\textsc{Find-Set}:] Nós chave por chamada: $O(\alpha(n))$;\\
        nós de bloco: custo amortizável ao potencial → $O(1)$ total por op.
    \end{description}
  \end{block}
  \vspace{0.3cm}
  \begin{block}{Resultado}
    Custo total $= O(m\,\alpha(n))$. No modelo de \emph{pointer machine} este limite é justo:\\
    Fredman--Saks (1989) provam $\Omega(m\,\alpha(n))$ no pior caso.
  \end{block}
\end{frame}

\begin{frame}{Tabela comparativa de complexidade}%% 58
  \begin{center}
  \renewcommand{\arraystretch}{1.4}
  \begin{tabular}{lll}
    \toprule
    \textbf{Implementação} & \textbf{Por op.\ (amort.)} & \textbf{$m$ operações} \\
    \midrule
    Lista ingênua               & $O(n)$         & $O(mn)$ \\
    Lista + union por peso      & $O(1)$ Find,\ amort.\ $O(\lg n)$ Union & $O(m+n\lg n)$ \\
    Floresta + rank             & $O(\lg n)$     & $O(m\lg n)$ \\
    Floresta + PC               & $O(\lg n)$     & $O(m\lg n)$ \\
    \textbf{Floresta + rank + PC} & $O(\alpha(n))$ & $\mathbf{O(m\,\alpha(n))}$ \\
    \bottomrule
    \multicolumn{3}{l}{\small $\alpha(n) \leq 4$ para qualquer $n$ prático} \\
  \end{tabular}
  \end{center}
\end{frame}

\begin{frame}{Recapitulando: Análise Amortizada}%% 59
  \begin{block}{O que aprendemos}
    \begin{itemize}
      \item Análise amortizada revela o custo médio real por operação
      \item Função de Ackermann cresce mais rápido que qualquer função primitiva recursiva
      \item $\alpha(n) \leq 4$ para qualquer $n$ concebível → prática é $O(1)$ amortizado
      \item Teorema 19.14 (CLRS): $O(m\,\alpha(n))$ no pior caso (justo no modelo de \emph{pointer machine})
    \end{itemize}
  \end{block}
\end{frame}

%% ================================================================
%%  SEÇÃO 8 — JAVA MODERNO
%% ================================================================
\section{Implementação em Java Moderno}

\begin{frame}{Java 17+: recursos relevantes}%% 60
  \begin{columns}[T]
    \begin{column}{0.48\textwidth}
      \begin{block}{Para Union-Find genérico}
        \begin{itemize}
          \item \textbf{Generics} (\texttt{<T>}): suportar qualquer tipo
          \item \textbf{HashMap}: mapeamento elemento → pai
          \item \textbf{Records}: imutabilidade para metadados
          \item \textbf{var}: inferência de tipo local
        \end{itemize}
      \end{block}
    \end{column}
    \begin{column}{0.48\textwidth}
      \begin{block}{Para uso em grafos}
        \begin{itemize}
          \item \textbf{Stream API}: processar arestas
          \item \textbf{Comparator}: ordenar por peso
          \item \textbf{instanceof pattern}: verificar tipos
          \item \textbf{sealed interfaces}: modelar tipos fechados
        \end{itemize}
      \end{block}
    \end{column}
  \end{columns}
\end{frame}

\begin{frame}[fragile,shrink=8]{Interface DisjointSet\textless T\textgreater{}}%% 61
\begin{lstlisting}
public interface DisjointSet<T> {
    /** Cria um conjunto unitario {x}. */
    void makeSet(T x);

    /** Retorna o representante do conjunto de x. */
    T find(T x);

    /**
     * Une os conjuntos de x e y.
     * @return true se estavam em conjuntos diferentes
     */
    boolean union(T x, T y);

    /** Verifica se x e y estao no mesmo conjunto. */
    default boolean connected(T x, T y) {
        return find(x).equals(find(y));
    }
}
\end{lstlisting}
\end{frame}

\begin{frame}[fragile,shrink=30]{Implementação com array de inteiros}%% 62
\begin{lstlisting}
public final class ArrayUnionFind {
    private final int[] parent;
    private final int[] rank;

    public ArrayUnionFind(int n) {
        parent = new int[n]; rank = new int[n];
        for (int i = 0; i < n; i++) parent[i] = i;
    }

    public int find(int x) {
        if (parent[x] != x) parent[x] = find(parent[x]);
        return parent[x];
    }

    public boolean union(int x, int y) {
        var rx = find(x); var ry = find(y);
        if (rx == ry) return false;
        if (rank[rx] < rank[ry]) { var t = rx; rx = ry; ry = t; }
        parent[ry] = rx;
        if (rank[rx] == rank[ry]) rank[rx]++;
        return true;
    }

    public boolean connected(int x, int y) { return find(x)==find(y); }
}
\end{lstlisting}
\end{frame}

\begin{frame}[fragile]{Record Java para aresta de grafo}%% 63
\begin{lstlisting}
// Record: imutavel, equals/hashCode/toString automaticos (Java 16+)
public record Edge<V>(V from, V to, double weight)
    implements Comparable<Edge<V>> {

    @Override
    public int compareTo(Edge<V> other) {
        return Double.compare(this.weight, other.weight);
    }
}

// Uso com var e generics:
var edges = List.of(
    new Edge<>("A", "B", 3.0),
    new Edge<>("B", "C", 1.5),
    new Edge<>("A", "C", 4.2)
);
\end{lstlisting}
\end{frame}

\begin{frame}[fragile,shrink=20]{Kruskal com Stream API}%% 64
\begin{lstlisting}
public static <V> List<Edge<V>> kruskal(
        Collection<V> vertices,
        Collection<Edge<V>> edges) {

    var uf = new GenericUnionFind<V>();
    vertices.forEach(uf::makeSet);

    return edges.stream()
        .sorted()                       // ordena por peso (Comparable)
        // efeito colateral seguro apenas em stream sequencial (nao use .parallel())
        .filter(e -> uf.union(e.from(), e.to())) // false = ciclo
        .collect(Collectors.toList());
}

// Uso:
var mst = kruskal(graph.vertices(), graph.edges());
mst.forEach(e -> System.out.printf("%s--%s: %.1f%n",
    e.from(), e.to(), e.weight()));
\end{lstlisting}
\end{frame}

\begin{frame}[fragile,shrink=15]{Detectar ciclos em grafo}%% 65
\begin{lstlisting}
public static <V> boolean hasCycle(
        Collection<V> vertices,
        Collection<Edge<V>> edges) {

    var uf = new GenericUnionFind<V>();
    vertices.forEach(uf::makeSet);

    for (var edge : edges) {
        // union retorna false se ja estao conectados = ciclo
        if (!uf.union(edge.from(), edge.to())) {
            return true;
        }
    }
    return false;
}
\end{lstlisting}
\vspace{0.2cm}
\begin{block}{Custo}
  $O(m\,\alpha(n))$ para $m$ arestas e $n$ vértices.
\end{block}
\end{frame}

\begin{frame}[fragile,shrink=5]{Testes JUnit 5 — Parte 1}%% 66
\begin{lstlisting}
import org.junit.jupiter.api.*;
import static org.junit.jupiter.api.Assertions.*;

class UnionFindTest {
    private ArrayUnionFind uf;

    @BeforeEach void setUp() { uf = new ArrayUnionFind(10); }

    @Test void singletons_are_not_connected() {
        assertFalse(uf.connected(0, 1));
        assertFalse(uf.connected(3, 7));
    }

    @Test void union_connects_elements() {
        uf.union(0, 1); uf.union(1, 2);
        assertTrue(uf.connected(0, 2));  // transitivo
        assertFalse(uf.connected(0, 3));
    }
\end{lstlisting}
\end{frame}

\begin{frame}[fragile,shrink=25]{Testes JUnit 5 — Parte 2}%% 67
\begin{lstlisting}
    @Test void find_is_idempotent() {
        uf.union(4, 5);
        int r1 = uf.find(4);
        int r2 = uf.find(4);
        assertEquals(r1, r2);
    }

    @Test void union_returns_false_if_already_connected() {
        uf.union(6, 7);
        boolean first  = uf.union(6, 7); // mesmos
        boolean second = uf.union(7, 6); // ordem invertida
        assertFalse(first);
        assertFalse(second);
    }

    @Test void kruskal_produces_mst_with_n_minus_1_edges() {
        var vertices = List.of(0, 1, 2, 3);
        var edges = List.of(new Edge<>(0,1,1.0), new Edge<>(1,2,2.0),
                            new Edge<>(2,3,3.0), new Edge<>(0,3,9.0));
        var mst = Kruskal.kruskal(vertices, edges);
        assertEquals(3, mst.size());
    }
}
\end{lstlisting}
\end{frame}

\begin{frame}{Recapitulando: Java Moderno}%% 68
  \begin{block}{O que aprendemos}
    \begin{itemize}
      \item Interface \texttt{DisjointSet<T>}: contrato limpo com \texttt{default}
      \item \texttt{GenericUnionFind<T>}: HashMap para tipos arbitrários
      \item \texttt{ArrayUnionFind}: arrays para performance máxima com inteiros
      \item Record \texttt{Edge<V>}: imutável, comparável, sem boilerplate
      \item Kruskal com Stream API: elegante e conciso
      \item JUnit 5: testes que documentam o comportamento esperado
    \end{itemize}
  \end{block}
\end{frame}

%% ================================================================
%%  SEÇÃO 9 — APLICAÇÕES
%% ================================================================
\section{Aplicações}

\begin{frame}{Kruskal: Árvore Geradora Mínima}%% 69
  \begin{center}
  \begin{tikzpicture}[
    nd/.style={circle,draw=UFSBlue,fill=UFSBlue!15,minimum size=0.65cm,font=\small\bfseries},
    mst/.style={very thick,UFSBlue},
    nst/.style={thick,dashed,UFSGray!50}
  ]
    \node[nd] (A) at (0,3)   {A};
    \node[nd] (B) at (2.5,4) {B};
    \node[nd] (C) at (5,3)   {C};
    \node[nd] (D) at (1.5,1.5){D};
    \node[nd] (E) at (3.5,1.5){E};

    %% MST edges (destacadas)
    \draw[mst] (A)--(B) node[midway,above left,font=\scriptsize] {2};
    \draw[mst] (B)--(C) node[midway,above right,font=\scriptsize] {3};
    \draw[mst] (A)--(D) node[midway,left,font=\scriptsize] {4};
    \draw[mst] (D)--(E) node[midway,below,font=\scriptsize] {1};

    %% Non-MST edges
    \draw[nst] (B)--(D) node[midway,fill=white,inner sep=1pt,font=\scriptsize] {8};
    \draw[nst] (C)--(E) node[midway,right,fill=white,inner sep=1pt,font=\scriptsize] {6};
    \draw[nst] (B)--(E) node[midway,fill=white,inner sep=1pt,font=\scriptsize] {7};
    \draw[nst] (A)--(E) node[midway,fill=white,inner sep=1pt,font=\scriptsize,below left] {9};

    \node[font=\scriptsize,UFSBlue] at (2.5,0.5)
      {MST: $\{D\text{-}E,A\text{-}B,B\text{-}C,A\text{-}D\}$, peso total = 10};
  \end{tikzpicture}
  \end{center}
\end{frame}

\begin{frame}{Kruskal: passo a passo com Union-Find}%% 70
  \begin{columns}[T]
    \begin{column}{0.55\textwidth}
      \begin{block}{Algoritmo de Kruskal}
        \begin{enumerate}
          \item Ordenar arestas por peso crescente
          \item Para cada aresta $(u,v)$ em ordem:\\
            \quad Se \textsc{Find-Set}$(u) \neq$ \textsc{Find-Set}$(v)$:\\
            \quad\quad adicionar à MST\\
            \quad\quad \textsc{Union}$(u,v)$
          \item Parar quando MST tem $n-1$ arestas
        \end{enumerate}
      \end{block}
    \end{column}
    \begin{column}{0.42\textwidth}
      \begin{exampleblock}{Execução (exemplo anterior)}
        \begin{tabular}{ll}
          $D$-$E$ (1): & \textsc{Union}(D,E) \\
          $A$-$B$ (2): & \textsc{Union}(A,B) \\
          $B$-$C$ (3): & \textsc{Union}(B,C) \\
          $A$-$D$ (4): & \textsc{Union}(A,D) \\
          $C$-$E$ (6): & ciclo — ignora \\
          $B$-$E$ (7): & ciclo — ignora \\
        \end{tabular}
      \end{exampleblock}
    \end{column}
  \end{columns}
\end{frame}

\begin{frame}[shrink=8]{Percolação}%% 71
  \begin{columns}[T]
    \begin{column}{0.47\textwidth}
      \begin{block}{Problema}
        Grade $n \times n$ de sítios.\\
        Cada sítio é aberto com probabilidade $p$.\\
        O sistema \textbf{percola} se existe caminho de sítios abertos do topo à base.
      \end{block}
      \vspace{0.3cm}
      \begin{block}{Solução com Union-Find}
        Nós virtuais: \textit{top} e \textit{bottom}.\\
        Ao abrir sítio $(i,j)$: unir com vizinhos abertos.\\
        Percola $\iff$ \textsc{Find-Set}(\textit{top}) = \textsc{Find-Set}(\textit{bottom}).
      \end{block}
    \end{column}
    \begin{column}{0.5\textwidth}
      \begin{tikzpicture}[scale=0.5]
        % Grid 5x5
        \foreach \i in {0,...,4} \foreach \j in {0,...,4}
          \draw[UFSGray!30] (\i,\j) rectangle (\i+1,\j+1);
        % Open cells (colored blue)
        \foreach \pos in {(0,4),(1,4),(2,4),(2,3),(2,2),(3,2),(3,1),(3,0),(4,0)}
          \fill[UFSBlue!40] \pos rectangle ++(1,1);
        % Path highlighted
        \draw[UFSRed,very thick,->]
          (0.5,4.5)--(1.5,4.5)--(2.5,4.5)--(2.5,3.5)--
          (2.5,2.5)--(3.5,2.5)--(3.5,1.5)--(3.5,0.5);
        \node[font=\scriptsize,UFSBlue] at (2.5,-0.6) {sistema percola};
      \end{tikzpicture}
    \end{column}
  \end{columns}
\end{frame}

\begin{frame}[shrink=5]{Outras aplicações}%% 72
  \begin{columns}[T]
    \begin{column}{0.48\textwidth}
      \begin{block}{Compiladores}
        \begin{itemize}
          \item Unificação de tipos em Prolog/ML
          \item Análise de equivalência de registradores
          \item Eliminação de variáveis duplicadas (SSA)
        \end{itemize}
      \end{block}
      \vspace{0.3cm}
      \begin{block}{Processamento de imagens}
        \begin{itemize}
          \item Segmentação de objetos em imagens binárias
          \item Flood fill com componentes conexas
        \end{itemize}
      \end{block}
    \end{column}
    \begin{column}{0.48\textwidth}
      \begin{block}{Redes e sistemas}
        \begin{itemize}
          \item Detecção de clusters em redes P2P
          \item Controle de congestionamento (grupos de enlace)
          \item Gerência de memória: blocos livres contíguos
        \end{itemize}
      \end{block}
      \vspace{0.3cm}
      \begin{block}{Jogos e simulações}
        \begin{itemize}
          \item Geração de labirintos (Kruskal)
          \item Física: simulação de materiais granulares
        \end{itemize}
      \end{block}
    \end{column}
  \end{columns}
\end{frame}

\begin{frame}{Resumo geral}%% 73
  \begin{center}
  \renewcommand{\arraystretch}{1.35}
  \begin{tabular}{lllll}
    \toprule
    \textbf{Técnica} & \textbf{Make} & \textbf{Find} & \textbf{Union} & \textbf{Total} \\
    \midrule
    Lista ingênua      & $O(1)$ & $O(1)$     & $O(n)$         & $O(mn)$ \\
    Lista + peso       & $O(1)$ & $O(1)$     & $O(\lg n)$     & $O(m{+}n\lg n)$ \\
    Floresta ingênua   & $O(1)$ & $O(n)$     & $O(1)$         & $O(mn)$ \\
    + rank             & $O(1)$ & $O(\lg n)$ & $O(1)$         & $O(m\lg n)$ \\
    + PC               & $O(1)$ & $O(\lg n)$ & $O(1)$         & $O(m\lg n)$ \\
    \textbf{+ rank+PC} & $O(1)$ & $O(\alpha)$& $O(1)$ & $\mathbf{O(m\alpha(n))}$ \\
    \bottomrule
  \end{tabular}
  \end{center}
  \vspace{0.2cm}
  \begin{center}
    \small$\alpha = \alpha(n) \leq 4$ para qualquer $n$ prático
  \end{center}
\end{frame}

\begin{frame}{Recapitulação final}%% 74
  \begin{block}{Jornada desta aula}
    \begin{enumerate}
      \item \textbf{Problema}: conectividade dinâmica — relação de equivalência
      \item \textbf{Lista encadeada}: Find $O(1)$, Union $O(n)$ → melhor: $O(m{+}n\lg n)$
      \item \textbf{Floresta}: Union $O(1)$, Find $O(\text{altura})$
      \item \textbf{Union by Rank}: controla altura → $O(\lg n)$ por Find
      \item \textbf{Path Compression}: achata o caminho → amortiza futuras Finds
      \item \textbf{Juntos}: $O(\alpha(n)) \approx O(4)$ — praticamente $O(1)$
      \item \textbf{Java}: GenericUnionFind<T>, Record, Stream, JUnit 5
    \end{enumerate}
  \end{block}
\end{frame}

\begin{frame}{Leituras recomendadas}%% 75
  \begin{itemize}
    \item \textbf{CLRS 4ª ed.}, Capítulo 19 — \textit{Data Structures for Disjoint Sets}:\\
      Seções 19.1–19.4: definições, listas, floresta, análise amortizada completa
    \item \textbf{Tarjan, R.E. (1975)}: \textit{Efficiency of a Good But Not Linear Set Union Algorithm} — artigo original da análise com $\alpha(n)$
    \item \textbf{Sedgewick \& Wayne} — \textit{Algorithms}, 4ª ed., Seção 1.5: abordagem prática com foco em Java
    \item \textbf{Skiena} — \textit{The Algorithm Design Manual}, Cap.\ 8: Union-Find no contexto de grafos
  \end{itemize}
\end{frame}

\begin{frame}{Exercícios propostos}%% 76
  \begin{enumerate}
    \item (CLRS 19.1-1) Se \textsc{Find-Set}$(x)$ for chamado em um nó $x$ cujo pai não é a raiz, o que acontece ao rank de $x$ após a compressão?
    \item Implemente \texttt{nComponents(int n, int[][] edges)} que conta o número de componentes conexas de um grafo não-dirigido.
    \item Mostre que a heurística de união por peso na lista encadeada garante $\lfloor\lg n\rfloor$ atualizações de ponteiro no máximo por elemento.
    \item Implemente percolação Monte Carlo: estime a probabilidade de percolação para $p \in [0,1]$ numa grade $20\times 20$ com 10\,000 experimentos.
    \item (Desafio) Implemente uma versão thread-safe de \texttt{ArrayUnionFind} usando \texttt{AtomicIntegerArray}.
  \end{enumerate}
\end{frame}

%% ================================================================
%%  REFERÊNCIAS
%% ================================================================
\section*{Referências}

\begin{frame}[allowframebreaks]{Referências}%% 79
  \nocite{*}
  \bibliographystyle{plain}
  \bibliography{referencias}
\end{frame}


\end{document}
